\documentclass[11pt,a4paper]{article}

% ── Packages ─────────────────────────────────────────────────────────────────
\usepackage[T1]{fontenc}
\usepackage[utf8]{inputenc}
\usepackage{mathptmx}          % Times for text and math
\usepackage{amsmath,amssymb,amsthm}
\usepackage{amsfonts}
% mathptmx drops \hbar — restore from AMSb
\AtBeginDocument{\let\hbar\relax
  \DeclareMathSymbol{\hbar}{\mathord}{AMSb}{"7E}}
\usepackage{geometry}
\usepackage{microtype}
\usepackage{hyperref}
\usepackage{booktabs}
\usepackage{enumitem}
\usepackage{setspace}
\usepackage{titlesec}
\usepackage{abstract}
\usepackage{natbib}

\geometry{
  top=1.1in, bottom=1.1in,
  left=1.15in, right=1.15in
}

% ── Section heading style (match RAQM) ───────────────────────────────────────
\titleformat{\section}{\normalfont\bfseries\large}{\thesection}{1em}{}
\titleformat{\subsection}{\normalfont\bfseries\normalsize}{\thesubsection}{1em}{}
\titleformat{\subsubsection}{\normalfont\bfseries\itshape\normalsize}{\thesubsubsection}{1em}{}
\titlespacing*{\section}{0pt}{14pt}{4pt}
\titlespacing*{\subsection}{0pt}{10pt}{3pt}
\titlespacing*{\subsubsection}{0pt}{8pt}{2pt}

% ── Theorem environments ──────────────────────────────────────────────────────
\theoremstyle{plain}
\newtheorem{theorem}{Theorem}[section]
\newtheorem{proposition}[theorem]{Proposition}
\newtheorem{corollary}[theorem]{Corollary}
\newtheorem{lemma}[theorem]{Lemma}
\theoremstyle{definition}
\newtheorem{definition}[theorem]{Definition}
\newtheorem{postulate}{Postulate}
\theoremstyle{remark}
\newtheorem{remark}[theorem]{Remark}
\newtheorem{claim}[theorem]{Claim}

% ── Abstract style ────────────────────────────────────────────────────────────
\renewcommand{\abstractnamefont}{\normalfont\bfseries}
\renewcommand{\abstracttextfont}{\normalfont\small}
\setlength{\absleftindent}{0.5cm}
\setlength{\absrightindent}{0.5cm}

% ── Hyperref ─────────────────────────────────────────────────────────────────
\hypersetup{
  colorlinks=true, linkcolor=black, citecolor=black, urlcolor=black,
  pdftitle={Relational Actualism Gravity and Cosmology (RAGC) v4},
  pdfauthor={Joshua F. Sandeman}
}

% ── Macros ───────────────────────────────────────────────────────────────────
\newcommand{\Pact}{\mathcal{P}_{\mathrm{act}}}
\newcommand{\Hon}{\mathcal{H}_{\mathrm{on}}}
\newcommand{\Sigmam}{\Sigma_{m}}
\newcommand{\Tact}{T^{\mu\nu}_{\mathrm{RA}}}
\newcommand{\Rmax}{R_{\mathrm{max}}}
\newcommand{\LLC}{\mathrm{LLC}}
\newcommand{\AP}{\mathrm{AP}}
\newcommand{\ket}[1]{|#1\rangle}
\newcommand{\bra}[1]{\langle #1|}
\newcommand{\braket}[2]{\langle #1 | #2 \rangle}

% ─────────────────────────────────────────────────────────────────────────────
\begin{document}

\begin{titlepage}
\centering
\vspace*{1cm}
{\LARGE\bfseries Relational Actualism Gravity and Cosmology (RAGC):\par}
\vspace{0.4cm}
{\large\itshape Covariance, The Local Ledger, Empirical Predictions,\\
and Mathematical Derivations in Actualism Gravity\par}
\vspace{1.2cm}
{\large Joshua F.\ Sandeman\par}
\vspace{0.3cm}
{\normalsize\itshape Independent Researcher, Salem, Oregon\par}
\vspace{0.3cm}
{\normalsize March 2026\par}
\vspace{0.3cm}
  {\small\textit{Preprint:}~\href{https://doi.org/10.5281/zenodo.19197999}
    {doi.org/10.5281/zenodo.19197999}\par}
\vspace{1.5cm}

\begin{abstract}
This paper formalizes the extension of Relational Actualism (RA) into the
non-inertial and curved spacetime regime (RAGC). We propose that the spacetime
manifold acts as an active, bandwidth-limited causal graph that drives the
objective pruning of the quantum state space. Modelling spacetime as a Directed
Acyclic Graph (DAG), we reconceive energy-momentum conservation as a strict
\emph{local ledger} balance condition at every actualization vertex. The paper
makes three empirical predictions that are testable independently of the open
mathematical derivations: (P1)~the inferred value of $H_0$ should correlate with
mean actualization density along the line of sight, explaining the Hubble Tension
without dark energy; (P2)~the vacuum energy contribution to the metric source term
is structurally zero, not fine-tuned, with implications for CMB power spectra
distinguishable from $\Lambda$CDM; and (P3)~daughter spacetimes nucleated at
causal severance boundaries inherit an effective cosmological constant determined
by the local actualization flux at nucleation---a dynamical alternative to the
anthropic landscape. Key mathematical results include: a spectral definition of the
Actualization Projector ($\Pact$) with a proof of vacuum energy suppression
(Theorem~\ref{thm:vacuum}); and the Ledger Partition Theorem
(Theorem~\ref{thm:cut}), establishing independent conservation across causal
severances, which has been \textbf{formally verified} in Lean~4/Mathlib.%
\footnote{The RA Lean~4 proof files (\texttt{RA\_AQFT\_Proofs\_v10.lean} and
\texttt{RA\_Proofs\_Lean4.lean}) will be made publicly available on arXiv
submission. The underlying Lean-QuantumInfo library (Meiburg, Lessa, and Soldati,
2025; arXiv:2510.08672) is publicly available at
\url{https://github.com/Timeroot/Lean-QuantumInfo}.}
Furthermore, the framework categorically prohibits causally silent dark matter
(WIMPs, axions, sterile neutrinos): a 100~GeV WIMP generates $10^{25}$--$10^{38}$
times less actualization flux than equivalent baryonic matter, falling below
the threshold for measurable gravitational lensing. Crucially, the observed
phenomenology attributed to dark matter --- flat rotation curves, Bullet Cluster
lensing offset, galaxy-galaxy lensing --- is accounted for without dark matter
via the $\nabla^2\!\ln\lambda$ mechanism and accumulated causal depth
$A_{\mathrm{RA}}$, classified across four distinct lensing regimes in the
companion paper RADM. The conjectural
mechanisms for flat galactic rotation curves and the Bullet Cluster offset---which
follow from the same actualization-density field---are developed in the companion
paper RADM. A new RA-native derivation of the
causal set generative measure is presented: the Benincasa-Dowker sign
criterion ($S_{\mathrm{BDG}}(v) > 0$) as the discrete actualization
filter yields a Poisson-CSG formula $c_k^{\mathrm{RA}}(x) \propto
e^{-\mu}\mu^k$ with $\mu = \lambda\,V_{\mathrm{coh}}\,p_{\mathrm{th}}$
that is local and environment-dependent.

\paragraph{Topological definition of $\lambda(x)$.}
The actualization density $\lambda(x)$ is defined \emph{purely combinatorially}
as $\lambda = N / \mathcal{V}$, where $N$ is the number of actualization
vertices in a local region and $\mathcal{V}$ is the volume of the corresponding
Alexandrov interval (causal diamond) $\mathcal{D}(p,q) = J^+(p) \cap J^-(q)$.
Crucially, $\mathcal{V}$ is measured by the \emph{longest directed chain}
(proper time) between the bounding past vertex $p$ and future vertex $q$ ---
not by any continuous background metric volume. The volume is emergent from the
causal order; the causal order does not live inside a pre-existing volume.
This definition is fully consistent with the Benincasa-Dowker continuum
limit~\citep{Benincasa2010} and makes $\lambda(x)$ a coordinate-independent,
graph-native quantity. There is no circularity: the metric is derived
from $\lambda$, not used to define it. A bandwidth-partition model
predicts $H_{\mathrm{local}} \approx 73.1\;\mathrm{km\,s^{-1}\,Mpc^{-1}}$
from the KBC void and a new falsifiable consequence: the Eridanus
supervoid ($\delta \approx -0.30$) should give $H \approx
75.9\;\mathrm{km\,s^{-1}\,Mpc^{-1}}$. The Jacobson coefficient
$\alpha = 2\pi$ is confirmed to within ${\sim}2\%$ when the full
baryonic disk mass is used --- consistent
with $\alpha = 2\pi$ exactly within the systematic uncertainty on
$v_{\mathrm{flat}}$ and $r_d$. Three mathematical derivations remain open and are precisely scoped for collaborators.
The fourth---proving that the Local Ledger Condition implies
$\nabla_\mu G^{\mu\nu} \equiv 0$ as the unique macroscopic field equation---is
closed by the companion paper \emph{Relational Actualism: The Einstein Field
Equations as the Unique Macroscopic Limit of the Actualization Graph}
(RACL), which proves $G_{\mu\nu} = 8\pi G\,
\mathcal{P}_{\mathrm{act}}[T_{\mu\nu}]$ with $\Lambda = 0$ as the unique
consistent macroscopic description of the RA-CSG, using the Lean-verified
LLC, the $\mathcal{P}_{\mathrm{act}}$ conservation theorem, and Lovelock's
uniqueness theorem.
The Einstein field equations are shown to be the unique consistent
macroscopic description of the RA-CSG by Lovelock's theorem, given
the Lean-verified conservation and covariance results; and the speed
of light is derived as the maximum causal bandwidth rate $c = l_P/t_P$,
with $E = \gamma mc^2$ following from discrete counting alone.

\medskip\noindent
\textbf{Keywords:} Causal Set Theory \textperiodcentered{} Actualization Projector
\textperiodcentered{} Local Ledger \textperiodcentered{} Vacuum Energy
\textperiodcentered{} Hubble Tension \textperiodcentered{} Graph Cut Theorem
\textperiodcentered{} Cosmological Constant \textperiodcentered{} Dark Matter \textperiodcentered{} Lean~4 Verification \textperiodcentered{} Benincasa-Dowker Action \textperiodcentered{} Poisson-CSG \textperiodcentered{} Eridanus Supervoid \textperiodcentered{} Lovelock Theorem \textperiodcentered{} Causal Bandwidth \textperiodcentered{} Lovelock Theorem \textperiodcentered{} Causal Bandwidth
\end{abstract}
\end{titlepage}

% ─────────────────────────────────────────────────────────────────────────────
\section{Introduction: The Limits of the Timeless Manifold}

The reconciliation of General Relativity and Quantum Mechanics remains the
central crisis of modern physics. However, this impasse may not merely be a
failure of specific models, but a failure of the underlying mathematical language.
Since the introduction of the Minkowski spacetime continuum, physics has relied
on a formalism that spatializes time. By treating time as a reversible parameter
$t$, standard physics enforces a ``block universe'' ontology---a timeless container
in which the future mathematically exists simultaneously with the past.

Crucially, standard Quantum Field Theory (QFT) suffers from this same mathematical
Platonism. Because QFT lacks a formal mathematical mechanism to represent the mere
fact that things happen, it fails to distinguish between mathematical potentia and
physical reality, granting full gravitational weight to unactualized virtual states.
The true crisis in physics is not merely that GR and QM are incompatible; it is
that both frameworks lack a mechanism for genuine becoming.

Relational Actualism (RA) proposes that time and causality are ontologically
primitive. Building upon relational frameworks \citep{Rovelli1996}, actualization
is not a subjective epistemic update; it is an objective, irreversible physical
interaction derived directly from on-shell boson exchanges and causal
severances \citep{Sandeman2026a}. Spacetime is not a passive background; it is an
active, relentlessly growing causal graph.

\paragraph{Time as the growing graph.}
A foundational ontological commitment of RAGC, distinct from its technical
content and worth stating explicitly before the formalism begins, is this:
\emph{time is not a dimension of a background manifold through which events are
located. Time is the process of the causal graph growing.}

In the block universe picture that special relativity appears to enforce,
past, present, and future co-exist as equally real coordinate slices. The
``flow'' of time is then either illusory or a feature of consciousness rather
than of physics. RA rejects this picture not on philosophical grounds but on
ontological ones: if actualization events are genuinely irreversible --- if the
vertices written into the causal DAG cannot be unwritten --- then the past is
fixed, the future is genuinely open, and the present is the growing frontier at
which possibilities become actualities. Temporal becoming is not a psychological
phenomenon; it is the physical process of the graph growing.

This identification has three immediate consequences. First, the arrow of time
requires no separate explanation. It is not a thermodynamic emergent phenomenon
explained by low-entropy initial conditions; it is encoded in the framework's most
primitive commitment, since actualization is by definition irreversible. Second,
the distinction between time and causality dissolves: to say that $A$ caused $B$
is to say that $A$'s actualization was a causal precursor to $B$'s in the DAG ---
there is no additional fact about ``temporal priority'' beyond causal precedence.
Third, the proper time of an observer is the count of actualization events along
their worldline, not a coordinate parameter. Time dilation --- both kinematic and
gravitational --- is the statement that different worldlines accumulate different
numbers of actualization events per coordinate second (Section~\ref{sec:timedilation}).

This paper proceeds as follows. Section~\ref{sec:postulates} states the postulates
of RAGC. Section~\ref{sec:engine_preview} previews the full generative algorithm. Section~\ref{sec:timedilation} derives relativistic time dilation in RA
terms. Section~\ref{sec:ledger} develops the local ledger formulation of
conservation. Section~\ref{sec:severances} analyzes causal severances and their
cosmological implications. Section~\ref{sec:predictions} states the empirical
predictions of RAGC explicitly, separated from the open derivation tasks.
Section~\ref{sec:math} presents the mathematical derivations, each developed
as far as currently tractable with hard walls identified precisely.

\paragraph{Note on formal verification.}
The Ledger Partition Theorem (Theorem~\ref{thm:cut}) and the Markov Blanket
Shielding Theorem have been formally verified in Lean~4 with Mathlib. The
verification establishes that the discrete graph-theoretic core of the Graph Cut
proof is logically watertight given the stated axioms. The Lean formalization is
\textbf{completely free of \texttt{sorry} tags} for the Graph Cut Theorem and its
supporting lemmas: the two sum-decomposition lemmas
(\texttt{sum\_outgoing\_decompose} and \texttt{sum\_incoming\_decompose}) have
been fully proved using \texttt{Finset} membership arguments and the
\texttt{arrow\_of\_time} condition.

A second Lean~4 file, \texttt{RA\_AQFT\_Proofs\_v10.lean}, formalises the
finite-dimensional AQFT proof program. It defines density matrices, unitary
conjugation, and the quantum relative entropy via the continuous functional
calculus (CFC), and proves \texttt{frame\_independence}
and \texttt{rindler\_stationarity} as compiled
theorems. The proof depends on three axioms: \texttt{vacuum\_lorentz\_invariant}
(the Poincar\'{e}-invariance of the Minkowski vacuum, a QFT physics input),
\texttt{Matrix.cfc\_conj\_unitary} (the CFC unitary-conjugation lemma, sourced
from the Lean-QuantumInfo library~\citep{Meiburg2025LQI}), and
\texttt{petz\_monotonicity} (the data processing inequality, proved in
Lean-QuantumInfo). All source files are available from the author.


% ─────────────────────────────────────────────────────────────────────────────
\section{The Engine of Becoming: A Preview of the Generative Algorithm}
\label{sec:engine_preview}
% ─────────────────────────────────────────────────────────────────────────────

The framework developed in this paper --- the Local Ledger Condition, the
Actualization Projector, the metric sourcing equation, the time dilation
account --- all follow from a single underlying algorithm that has not yet
been stated in its most complete form. The full treatment is the subject of
the companion paper \emph{The Engine of Becoming}~\citep{Sandeman2026RAEB}, but
a preview is essential here because it reveals that RAGC's individual results
are not separate claims but consequences of one unified generative process.

The algorithm is a \textbf{Stochastic Graph Rewriting System}. It replaces
both the Schr\"{o}dinger equation and the Einstein field equations with a
single iterative procedure operating on the causal DAG. It is the precise
formal answer to the question: what is the universe \emph{doing}, step by
step, as the causal graph grows?

\paragraph{The five steps.}
Let the universe at step $n$ be defined by the graph $G_n = (V_n, E_n)$,
where $V_n$ is the set of actualized vertices and $E_n$ is the set of directed
causal edges.

\begin{description}[leftmargin=1.5em, itemsep=4pt]
  \item[\textbf{Step 1 --- The Open Future.}]
  The locus of potential physical interaction is the boundary of the graph:
  the set of outgoing edges that have not yet terminated at a subsequent vertex.
  Only events connected to this boundary can be actualized next. The past is
  fixed; only the frontier is live.

  \item[\textbf{Step 2 --- The Quantum Possibility Space.}]
  From the boundary $\partial G_n$, the system explores the space of
  unactualized candidate interactions. Transition amplitudes are computed
  via the standard Feynman path integral over the effective background metric
  $g^{(n)}_{\mu\nu}$ generated by the current graph state. These candidates are
  virtual: they do not belong to $V_n$ and do not source gravitational curvature.
  This is why the vacuum does not gravitate in RA --- unactualized potentia
  carry no metric weight.

  \item[\textbf{Step 3 --- The Kinematic Snap.}]
  Not all candidates can become permanent facts. A candidate vertex $v_*$ is
  promoted to the valid actualization set $\mathcal{V}_{\mathrm{valid}}$ only
  if it satisfies the actualization criterion: irreversible increase in
  entanglement entropy with the asymptotic environment. In the perturbative
  regime this is implemented by the on-shell threshold
  $p^{\mu}p_{\mu} \geq m^2c^2$. This is the objective, observer-independent
  collapse mechanism --- not an axiom but a consequence of the process ontology.

  \item[\textbf{Step 4 --- The Stochastic Update.}]
  The universe advances by selecting exactly one vertex from
  $\mathcal{V}_{\mathrm{valid}}$, weighted by the Born rule measure
  $W(v_*) = |\mathcal{A}(v_*)|^2$:
  \begin{equation}
    P(v_{n+1}) = \frac{W(v_{n+1})}{\displaystyle\sum_{v_i \in \mathcal{V}_{\mathrm{valid}}} W(v_i)}.
    \label{eq:stochastic_update}
  \end{equation}
  The graph updates: $V_{n+1} = V_n \cup \{v_{n+1}\}$. All contradictory
  virtual histories branching from the same causal antecedents are pruned.
  This is wave function collapse --- not as a postulate, but as the
  inevitable consequence of a single irreversible addition to the causal ledger.

  \item[\textbf{Step 5 --- The Hydrodynamic Limit (Emergent GR).}]
  Gravity is not mediated by a particle during Step~2. It is the large-scale
  statistical geometry of the actualized graph, recalculated \emph{only after}
  Step~4 from the updated vertex set. The local actualization density
  $\lambda(x)$ is computed over a mesoscopic coarse-graining volume, and the
  effective spacetime metric updates via the RA-modified field equations
  (Section~\ref{sec:pact}). This new metric then serves as the background
  for Step~2 of the next iteration. Gravity and quantum mechanics operate on
  different steps of the same algorithm --- which is precisely why they have
  been so resistant to unification by frameworks that try to run them
  simultaneously.
\end{description}

\paragraph{The BMV prediction.}
Because the spacetime metric is recalculated strictly in Step~5 --- predicated
entirely on the set of actualized vertices $V_n$ --- it is formally impossible
for a quantum superposition (Step~2) to source a superposed gravitational
field. Relational Actualism therefore predicts a strict null result for
experiments attempting to observe gravity-mediated entanglement between
superposed masses (the Bose-Marletto-Vedral protocol~\citep{Bose2017,
Marletto2017}). Spacetime cannot be caught in a superposition because
spacetime is the accumulated ledger of interactions that have already
irreversibly snapped. This is the most near-term discriminating prediction
of the full Engine of Becoming framework.

\paragraph{How RAGC's results fit the algorithm.}
Every result in this paper is a consequence of one or more steps of this
algorithm. The Local Ledger Condition (Section~\ref{sec:ledger}) is the
conservation law enforced by the Step~4 pruning. The Actualization Projector
$\Pact$ (Section~\ref{sec:pact}) is the formal implementation of the
Step~3 threshold. The metric sourcing equation
(Derivation~5, Section~\ref{sec:d5}) is the Step~5 hydrodynamic limit.
The vacuum energy suppression (Theorem~\ref{thm:vacuum}) follows from the
Step~2 non-sourcing of virtual states. The Hubble tension prediction
(Prediction~1) arises from the spatial variation in the Step~5 coarse-graining
rate. RAGC is not a collection of separate results; it is the first systematic
derivation of the consequences of a single algorithm.

% ─────────────────────────────────────────────────────────────────────────────
\section{Postulates and Definitions}
\label{sec:postulates}

\begin{postulate}[Ontic State]
The fundamental state of a physical system is described by its density operator
$\rho$ on a separable Hilbert space $\mathcal{H}$.
\end{postulate}

\begin{postulate}[Actualization]
Actualization is the objective, irreversible suppression of off-diagonal coherences
in $\rho$, driven by environmental entanglement. This criterion runs parallel to
dynamical collapse models \citep{Bassi2013}, with the key distinction that the
trigger in RA is a specific on-shell boson exchange condition rather than a
stochastic gravitational threshold. The actualization criterion---irreversible
increase in entanglement entropy with the asymptotic environment---is
\emph{sufficient} for actualization in all regimes. Its necessity in full
generality, and its extension to gauge theories with non-trivial edge mode
structure, are open problems inherited from RAQM \citep{Sandeman2026a}.
\end{postulate}

\begin{postulate}[The Environment]
The environment driving actualization is the coarse-grained geometry of the
spacetime manifold itself.
\end{postulate}

\begin{postulate}[Bandwidth Limit]
The manifold possesses a fundamental information-theoretic channel capacity.
There exists a maximum rate $\Rmax$ at which the local spacetime geometry can
process quantum state transitions.
\end{postulate}

\begin{postulate}[Actualization Projector]
The macroscopic metric is sourced exclusively by on-shell interactions. Off-shell
virtual fluctuations are mathematically unactualized potentia. We formally propose
a projector $\Pact$ that restricts the quantum stress-energy tensor to states
satisfying $p^{\mu}p_{\mu} = m^{2}c^{2}$. The precise definition of this
projector, and the proof that it suppresses the vacuum energy contribution, is
developed in Section~\ref{sec:pact}.
\end{postulate}

\begin{remark}
Unlike Postulates~1--4, which state ontological commitments, Postulate~5
constitutes a conjecture whose formal elaboration is an open mathematical problem.
We state it here to motivate the framework; the mathematical development in
Section~\ref{sec:pact} constitutes the first rigorous partial derivation.
\end{remark}


% ─────────────────────────────────────────────────────────────────────────────
\section{Proper Time and Relativistic Dilation in the Actualization Framework}
\label{sec:timedilation}
% ─────────────────────────────────────────────────────────────────────────────

\subsection{Proper Time as Actualization Count}

The identification of time with the growing causal graph has a precise
corollary for the meaning of proper time. In standard GR, proper time is
the spacetime interval along a worldline:
\begin{equation}
  d\tau^2 = g_{\mu\nu}\,dx^\mu\,dx^\nu.
  \label{eq:proper_time_gr}
\end{equation}
In RA, the spacetime metric is emergent from the causal graph rather than
fundamental. The primitive quantity is not the interval but the count of
actualization events along the worldline. We define:

\begin{definition}[Proper Actualization Count]
The \emph{proper actualization count} $\mathcal{N}(\gamma)$ of a worldline
$\gamma$ is the number of on-shell actualization events that occur along $\gamma$
between two events $p, q \in \mathcal{G}$.
\end{definition}

The proper time measured by a physical clock is $\tau = \mathcal{N}(\gamma)\cdot t_P$:
a clock is a physical system that counts its own actualization events, and
proper time is that count in Planck units. This identification is exact at
every scale including the Planck scale --- it is not a continuum limit.
In the macroscopic regime $\mathcal{N}(\gamma) \gg 1$, the discrete count
converges to the standard GR proper time integral $\int d\tau$, but the
RA definition is the integer count itself, prior to any continuum
approximation.

\subsection{Special-Relativistic Time Dilation}
\label{subsec:sr_dilation}

Consider two observers $\mathcal{O}$ and $\mathcal{O}'$ in relative motion, with
$\mathcal{O}'$ moving at velocity $v$ relative to $\mathcal{O}$. In the RA
account, each observer accumulates actualization events at a rate determined by
the on-shell interactions along their respective worldlines.

The actualization criterion is Lorentz-invariant: whether a given virtual exchange
constitutes an actualization event is determined by the scalar condition
$p^\mu p_\mu = m^2 c^2$, which is unchanged by any Lorentz transformation.
However, the \emph{rate} at which such events accumulate along a given worldline
is frame-dependent, because the Lorentz boost changes the energy-momentum of
mediating bosons:
\begin{equation}
  p^\mu \;\longrightarrow\; \Lambda^\mu{}_\nu\, p^\nu,
  \label{eq:lorentz_boost}
\end{equation}
where $\Lambda$ is the Lorentz transformation. In the rest frame of $\mathcal{O}'$,
the typical boson energy is $\langle E \rangle$. In the frame of $\mathcal{O}$,
the same bosons are blue- or red-shifted; the threshold crossing rate --- and
therefore the actualization rate per coordinate second --- changes accordingly.

The ratio of proper actualization counts for the two observers between the same
pair of events is:
\begin{equation}
  \frac{\mathcal{N}(\gamma')}{\mathcal{N}(\gamma)}
  \;=\; \sqrt{1 - \frac{v^2}{c^2}}
  \;=\; \frac{1}{\gamma_{\mathrm{LF}}},
  \label{eq:sr_dilation}
\end{equation}
where $\gamma_{\mathrm{LF}} = (1-v^2/c^2)^{-1/2}$ is the Lorentz factor.
This is the standard time dilation formula, now derived from the actualization
count rather than from the metric interval. The moving clock accumulates fewer
actualization events per coordinate second not because time ``flows more slowly''
in some absolute sense, but because its worldline through the causal graph
traverses a different density of on-shell threshold crossings.

\begin{remark}
This account makes the origin of SR time dilation transparent in a way that the
geometric account does not. The geometric account says: the spacetime interval is
shorter along the moving worldline. The RA account says: the moving observer
undergoes fewer irreversible causal interactions per coordinate second. The
second is more fundamental, because it gives a physical mechanism (actualization
rate) for what the first describes geometrically. The two accounts agree quantitatively in the macroscopic regime
$\mathcal{N} \gg 1$; the RA account is more primitive and exact at all scales,
including the Planck scale where the continuum interval is undefined.
\end{remark}

% ─── Bandwidth subsection ───────────────────────────────────────────────────
\subsection{The Speed of Light as Causal Bandwidth Ceiling}
\label{subsec:bandwidth}

The discrete-first perspective yields a derivation of $c$ as a fundamental
limit and of the relativistic energy relation $E = \gamma mc^2$ from RA
primitives alone, without importing either result from continuum physics.

\paragraph{Rest-frame actualization rate.}
A particle of mass $m$ at rest must maintain its on-shell identity via
actualization events. From the on-shell condition $p^\mu p_\mu = m^2c^2$
and the RA action quantization (each Kinematic Snap carries action $h$):
\begin{equation}
  f_0 = \frac{mc^2}{h},
  \label{eq:rest_freq}
\end{equation}
the Compton frequency, here derived from the actualization criterion rather
than imported from non-relativistic quantum mechanics.

\paragraph{Boost as bandwidth redistribution.}
The four-momentum $p^\mu = \hbar k^\mu$ is the actualization current:
$p^0 = E/c$ is the temporal actualization rate and $p^i$ is the spatial
causal connection rate. The on-shell invariant $p^\mu p_\mu = m^2c^2$ is
the statement that the total actualization current magnitude is frame-independent.
Boosting to velocity $v$, the temporal component transforms as
$f_{\mathrm{coord}} = \gamma f_0$, giving immediately:
\begin{equation}
  E = h\,f_{\mathrm{coord}} = \gamma m c^2.
  \label{eq:rel_energy_derived}
\end{equation}
This derivation uses only the on-shell condition, action quantization per
Kinematic Snap, and the Lorentz transformation of the frequency four-vector.

\paragraph{Why $c$ is the ceiling.}
A massless particle ($m = 0$) has $f_0 = 0$: no internal bandwidth is
consumed by mass-maintaining actualization. All causal bandwidth goes to
spatial propagation, giving speed exactly $c = l_P/t_P$ --- the ratio of
the two fundamental graph scales. A massive particle approaching $c$ requires
$\gamma \to \infty$ temporal bandwidth to maintain its on-shell identity
while propagating. Since the causal graph has maximum vertex rate
$c/l_P$ per Planck volume, this is unachievable for finite energy:
\begin{equation}
  c = \frac{l_P}{t_P}
  \quad \text{is the maximum causal propagation rate of the RA graph.}
  \label{eq:c_ceiling}
\end{equation}
The speed of light is the ratio of the two fundamental graph scales.
Its role as a universal speed limit is the statement that no massive
particle can devote its entire causal bandwidth to spatial propagation.

\paragraph{Event horizons as local bandwidth saturation.}
The Causal Severance of an event horizon is the same mechanism at the
scale of the graph boundary: when the local graph's expansion rate reaches
$c$, all local causal bandwidth is committed to boundary dynamics and no
outgoing edges remain. The daughter spacetime nucleation
(Section~\ref{sec:severances}) is the release of this saturated
bandwidth into a new causal graph --- the $v \to c$ limit applied to
the graph boundary rather than to a worldline within it.

\subsection{Gravitational Time Dilation}
\label{subsec:grav_dilation}

Gravitational time dilation in RA follows from the metric sourcing equation
rather than from the Lorentz boost. A region of high actualization density
$\lambda(x)$ corresponds to a densely woven causal graph --- many vertices per
unit coordinate volume, tightly constrained by existing causal relations. A
clock embedded in such a region is embedded in a dense causal structure; each
new actualization event it undergoes is constrained by many prior events.

In the weak-field limit, the gravitational potential $\Phi$ is sourced by
$\lambda(x)$ via the RA Poisson equation:
\begin{equation}
  \nabla^2\Phi = 4\pi G\,\rho_A(x),
  \qquad
  \rho_A \;\propto\; \frac{A_{\mathrm{RA}}(x)}{V},
  \label{eq:poisson_grav_dilation}
\end{equation}
where $\rho_A$ is the accumulated causal depth density. A clock at potential
$\Phi$ accumulates actualization events at a rate modified by the local graph
topology. In the Newtonian limit:
\begin{equation}
  \frac{\mathcal{N}(\gamma_{\Phi})}{\mathcal{N}(\gamma_0)}
  \;\approx\; 1 + \frac{\Phi}{c^2},
  \label{eq:grav_dilation}
\end{equation}
where $\gamma_0$ is a reference worldline far from the source and $\gamma_\Phi$
is a worldline at potential $\Phi < 0$. Clocks deep in a gravitational well
accumulate fewer actualization events per coordinate second --- they run slower
--- because the dense causal graph in that region constrains the rate at which
new on-shell events can be added to the local subgraph.

\paragraph{Unification.}
Both kinematic and gravitational time dilation are, in RA, the same phenomenon:
the rate of accumulation of irreversible causal events along a worldline. The
Lorentz boost changes this rate by transforming boson four-momenta; the
gravitational field changes it by modifying the local causal graph density.
Special-relativistic and general-relativistic time dilation are two regimes of
the same underlying actualization-rate calculation.

\paragraph{Hard wall.}
The formal derivation of equations~\eqref{eq:sr_dilation}
and~\eqref{eq:grav_dilation} from the covariant RA field equations requires
establishing Bianchi compatibility ($\nabla_\mu \langle T^{\mu\nu}_{\mathrm{RA}}
\rangle = 0$) on general curved backgrounds --- Derivation~3, hard wall~iv.
The qualitative account above is derivable from the actualization criterion and
the proper-time identification; the precise numerical coefficients require the
full covariant framework. What is not conjectural is the identification of proper
time with the actualization count: this follows directly from the ontological
commitment that time is the growing graph.

% ─────────────────────────────────────────────────────────────────────────────
\section{The Local Ledger: Conservation on a Growing Graph}
\label{sec:ledger}

In RA, the universe is an actively growing causal graph. Extending the topology
of Causal Set Theory \citep{Dowker2005} and its classical sequential growth
dynamics \citep{RideoutSorkin2000}, new actualization events continuously add
vertices and directed edges. Conservation laws in RAGC must be reconceived as
local ledger properties rather than global constants.

Define the actualization graph $G = (V, E)$, where $V$ is the set of actualization
events and $E$ is the set of directed causal edges. Each vertex $v \in V$ carries
a charge vector $\mathbf{q}(v) \in \mathbb{R}^{n}$, whose components are the
conserved charges (energy-momentum, electric charge, baryon number, etc.).

\begin{definition}[Local Ledger Condition (LLC)]
For every vertex $v \in V$:
\begin{equation}
  \sum_{\substack{u \in V \\ (u,v)\in E}} \mathbf{q}(u)
  \;=\;
  \sum_{\substack{w \in V \\ (v,w)\in E}} \mathbf{q}(w).
\end{equation}
\end{definition}

That is, the sum of incoming charges equals the sum of outgoing charges at every
vertex. In the continuum limit of a dense graph (vertex density $\lambda \to
\infty$), the LLC converges to the standard Noether continuity equation
$\nabla_{\mu}J^{\mu} = 0$. A sketch of this convergence is given in
Section~\ref{sec:continuum}.

% ─────────────────────────────────────────────────────────────────────────────
\section{Causal Severances and Cosmological Expansion}
\label{sec:severances}

\subsection{Event Horizons and Ledger Partitioning}

In RAGC, an event horizon acts as an absolute actualization boundary,
topologically severing the causal graph. The Graph Cut Theorem
(Theorem~\ref{thm:cut}) establishes rigorously that the local conservation
condition holds independently on each side of the horizon cut.

\paragraph{Scope clarification.}
This result resolves the \emph{energy bookkeeping} version of the Hawking
radiation problem---specifically, the question of whether total energy is conserved
across the horizon. In RAGC, the question ``is $Q_{\mathrm{interior}} +
Q_{\mathrm{exterior}}$ conserved?'' is not well-posed: after the causal cut,
$Q_{\mathrm{interior}}$ is inaccessible to the exterior subgraph by definition of
an event horizon. The claim is \emph{not} that the full unitarity question (whether
quantum information is preserved) is resolved---that requires the complete quantum
gravity treatment of $\Pact$ in the non-perturbative regime, which is an explicit
open problem (Section~\ref{sec:hardwall}). This distinction is maintained
throughout.

When the actualization flux approaches the bandwidth limit $\Rmax$, a topological
severance occurs, nucleating a causally orthogonal region. This is not merely an
analogy with black hole formation: it is a precise claim about what happens when
the local rate of on-shell interactions saturates the information-processing
capacity of the spacetime manifold. The DAG can no longer weave new edges into the
existing causal structure fast enough to maintain connectivity; a new, topologically
disconnected subgraph is nucleated.

The ongoing accretion from the parent spacetime feeds an external energy flux into
this daughter region through the severance boundary. In RAGC, this flux is not
arbitrary --- it is governed by the Local Ledger Condition at the boundary vertices
(Theorem~\ref{thm:cut}), which constrains the charge flow across the cut. The
energy density that the daughter spacetime inherits at nucleation depends on the
actualization flux at the moment of severance, which is determined by the local
matter density and interaction rates in the parent spacetime at that moment.

This provides a candidate mechanism for cosmological constant selection: different
daughter spacetimes, nucleated at different moments and locations in the parent
DAG, inherit different initial energy densities from the boundary flux. The
effective cosmological constant of each daughter spacetime is thus not a
fundamental free parameter but a derived quantity --- selected by the actualization
conditions at the moment of causal severance. Spacetimes nucleated from high-flux
regions inherit higher initial energy densities; those from low-flux regions
inherit lower ones.

\paragraph{Relationship to the anthropic landscape.}
This mechanism offers a physically grounded alternative to the string theory
landscape as an account of cosmological constant selection. In the landscape
picture, the cosmological constant is selected anthropically from an enormous
ensemble of vacua. In RAGC, the selection mechanism is dynamical and local: the
value of $\Lambda$ in a daughter spacetime is determined by the actualization flux
at its nucleation event, not by a random draw from a distribution of vacua. This
is a sharper prediction in the sense that it ties the cosmological constant to
measurable properties of the parent spacetime's matter distribution and interaction
rates.

\emph{We note this entire discussion as a conjecture, not a derivation.} The
quantitative prediction --- establishing the functional relationship between the
boundary flux at nucleation and the effective $\Lambda$ of the daughter spacetime
--- requires the non-perturbative treatment of $\Pact$ (Section~\ref{sec:hardwall})
and the resolution of the Bianchi compatibility problem. It is, however, a
\emph{testable} conjecture in principle: if our universe is a daughter spacetime
in this sense, its cosmological constant should reflect the actualization flux
conditions at the nucleation event of our causal structure, which in turn should
leave an imprint in the large-scale structure of the early universe. Identifying
that imprint is a long-range observational target for the program.

\subsection{Intrinsic Graph Growth and the Hubble Tension}

Independent of extrinsic feeding, the universe experiences intrinsic graph growth.
As matter and energy interact, the continuous actualization of on-shell events
physically weaves new topological edges. Actualization-dense regions (galaxy
clusters) have high edge density; their metric is, in the continuum limit, tightly
constrained. Intergalactic voids, with low actualization rate density $\lambda(x)$,
permit maximum intrinsic topological expansion.

The mechanism rests on a distinction between the \emph{rate of vertex addition} and
the \emph{coordinate volume subtended per vertex}. In dense regions, vertices are
added rapidly but each is tightly coupled to neighbors, leaving little room for the
metric to expand between events. In voids, vertices are added rarely, but each
subtends a large coordinate volume because no neighboring interactions constrain
the local geometry. The metric breathes freely between sparse causal stitches in a
way it cannot in a densely woven region.

This generates the prediction $H_{\mathrm{eff}}(x) \propto 1/\lambda(x)$, which
provides the mechanism for the Hubble Tension. The full statement of this
prediction, its distinguishing observational signature, and its distinction from
$\Lambda$CDM and early dark energy models are given in
Section~\ref{sec:predictions}, Prediction~1. The quantitative derivation requires
fixing the generative measure of the DAG (Section~\ref{sec:algebra}).

% ─────────────────────────────────────────────────────────────────────────────
\section{Empirical Predictions}
\label{sec:predictions}

The framework makes predictions that are testable observationally and experimentally,
independently of whether every mathematical derivation in Section~\ref{sec:math}
is complete. This section states them plainly. A theory earns its standing by
making predictions, not by deferring all claims until every derivation is polished.
The quantitative precision of each prediction depends on completing the derivations
identified in Section~\ref{sec:math}; the qualitative direction and distinguishing
signature of each prediction follows from the core postulates alone.

\subsection{Prediction 1: The Hubble Tension as a Density-Dependent Expansion Rate}

RAGC predicts that the effective cosmic expansion rate $H_{\mathrm{eff}}(x)$ is
inversely proportional to the local actualization density $\lambda(x)$:
\begin{equation}
  H_{\mathrm{eff}}(x) \;\propto\; \frac{1}{\lambda(x)}.
  \label{eq:hubble_pred}
\end{equation}
This is a direct structural consequence of the DAG metric: in regions of high
actualization density (galaxy clusters, filaments), the metric is tightly woven
and resists expansion; in regions of low actualization density (intergalactic
voids), the metric grows freely between sparse causal events.

\paragraph{The Hubble Tension.}
Early-universe probes (CMB temperature anisotropies, baryon acoustic oscillations)
sample the mean actualization rate of the dense, radiation-dominated early universe.
Late-universe probes (Cepheid distance ladders, Type~Ia supernovae, strong lensing
time delays) are dominated by lines of sight through low-$\lambda$ void-dominated
regions. RAGC predicts these two probe families measure genuinely different physical
quantities: $H_{\mathrm{eff}}$ in different density regimes. The persistent
$\sim\!4$--$6\sigma$ discrepancy between early and late measurements is therefore
not a systematic error but a real signal of density-dependent expansion.

\paragraph{Distinguishing signature.}
This prediction is observationally discriminating in the following sense: the
inferred $H_0$ from any probe should correlate with the mean matter density along
its line of sight. Void-dominated survey geometries should return systematically
higher $H_0$; cluster-dominated geometries should return lower values. This
line-of-sight dependence is \emph{not} predicted by $\Lambda$CDM (which treats
$H_0$ as a global constant) nor by competing early dark energy models (which shift
the sound horizon uniformly). RAGC predicts a specific functional relationship
between $H_0$ and the integrated mass density along each line of sight,
quantifiable once the generative measure of the DAG is fixed (Section~\ref{sec:algebra}).

\paragraph{No dark energy required.}
The differential expansion rate between dense and sparse regions is a structural
consequence of the actualization ontology itself, requiring no additional dark
energy ingredient.

\subsection{Prediction 1b: The Bandwidth-Partition Formula and the Eridanus Test}
\label{sec:hubble_partition}

The mechanism of Prediction~1 can be made quantitative without requiring the
full derivation of the generative measure $c_k$. The total actualization
bandwidth is partitioned between matter-binding edges and spatial-expansion
edges:
\begin{equation}
  \lambda_{\mathrm{total}} = \lambda_m + \lambda_{\mathrm{exp}},
  \label{eq:bandwidth_partition}
\end{equation}
where $\lambda_m$ is the rate of causal connections that bind matter
(internal knots of the causal graph) and $\lambda_{\mathrm{exp}}$ is the
rate of connections that grow the graph spatially (transverse expansion edges).
Because the total bandwidth is locally conserved, a deficit in matter density
frees bandwidth for expansion. For a region with fractional density contrast
$\delta \equiv (\rho - \bar\rho)/\bar\rho$ relative to the cosmic mean, the
local expansion bandwidth is:
\begin{equation}
  \lambda_{\mathrm{exp}}^{\mathrm{local}}
  = \bar\lambda_{\mathrm{exp}} - \delta\,\bar\lambda_m,
\end{equation}
giving the ratio of local to global Hubble rates:
\begin{equation}
  \frac{H_{\mathrm{local}}}{\bar H}
  = 1 - \delta\,\frac{\bar\lambda_m}{\bar\lambda_{\mathrm{exp}}}.
  \label{eq:hubble_partition_full}
\end{equation}

\paragraph{One-parameter fit to the KBC void.}
The KBC void~\citep{Kenworthy2019} is a well-documented underdensity of radius
${\sim}300\;\mathrm{Mpc}$ with density contrast $\delta \approx -0.20$.
Using the SH0ES local measurement $H_{\mathrm{local}} = 73.04 \pm
1.04\;\mathrm{km\,s^{-1}\,Mpc^{-1}}$~\citep{Riess2022} and the Planck CMB
value $\bar H = 67.4\;\mathrm{km\,s^{-1}\,Mpc^{-1}}$, the implied ratio is:
\begin{equation}
  \frac{\bar\lambda_m}{\bar\lambda_{\mathrm{exp}}}
  = \frac{H_{\mathrm{local}}/\bar H - 1}{-\delta}
  = \frac{0.084}{0.20} \approx 0.42.
  \label{eq:lambda_ratio_fit}
\end{equation}
This is a one-parameter fit. The ratio $\bar\lambda_m/\bar\lambda_{\mathrm{exp}}$
will be derived from first principles once the BDG-filter MCMC
(Section~\ref{sec:algebra}, Hard Wall~1) is complete; until then it is
observationally calibrated from the KBC void.

\paragraph{Falsifiable prediction: the Eridanus supervoid.}
The Eridanus supervoid has an estimated density contrast $\delta \approx -0.30$
and diameter ${\sim}300\;\mathrm{Mpc}$~\citep{Naidoo2016}. Substituting into
equation~\eqref{eq:hubble_partition_full} with the fitted ratio:
\begin{equation}
  H_{\mathrm{Eridanus}}
  = 67.4 \times \bigl(1 - (-0.30)(0.42)\bigr)
  \approx 75.9\;\mathrm{km\,s^{-1}\,Mpc^{-1}}.
  \label{eq:eridanus_prediction}
\end{equation}
This prediction is RA-native and structurally distinct from $\Lambda$CDM (which
predicts no dependence of $H_0$ on the density of the observing environment) and
from standard void models (which produce at most ${\sim}1$--$2\;\mathrm{km\,s^{-1}\,Mpc^{-1}}$
corrections from peculiar velocities alone). A measurement of $H$ in the
direction of the Eridanus supervoid consistent with ${\sim}75.9\;\mathrm{km\,s^{-1}\,Mpc^{-1}}$
would provide strong support for the bandwidth-partition mechanism; a value
near $67$--$68\;\mathrm{km\,s^{-1}\,Mpc^{-1}}$ would falsify it.

\paragraph{Linear dependence on $\delta$: the full observational test.}
Equation~\eqref{eq:hubble_partition_full} predicts that $H_0$ measurements
should vary \emph{linearly} with the mean density contrast $\bar\delta$ along
each line of sight, with slope $-\bar H \cdot (\bar\lambda_m/\bar\lambda_{\mathrm{exp}})
\approx -28.3\;\mathrm{km\,s^{-1}\,Mpc^{-1}}$ per unit $\delta$. This is a
falsifiable functional form, not merely a direction. Current data from the
Carnegie-Chicago Hubble Program, H0LiCOW, and TDCOSMO, stratified by
line-of-sight density, already contain sufficient statistical power to test
this slope prediction against the $\Lambda$CDM null of zero slope.

\subsection{Prediction 2: Vacuum Energy Suppression and CMB Implications}

RAGC predicts that the vacuum energy contribution to the metric source term is
structurally zero:
\begin{equation}
  \langle 0 \,|\, \Tact \,|\, 0 \rangle = 0,
\end{equation}
as proved in Theorem~\ref{thm:vacuum}. At the ontological level,
this is a structural consequence of the actualization postulate: the metric is
sourced only by actualized (on-shell) states, and the vacuum contains none.
At the operational level in semiclassical gravity, this motivates the
vacuum-subtraction prescription (equation~\eqref{eq:ra_prescription} in
Section~\ref{sec:pact}), which implements the same suppression covariantly
within the renormalized framework. The two framings are not in tension:
the ontological claim explains \emph{why} the vacuum does not gravitate in RA;
the prescription specifies \emph{how} to compute this consistently in the
presence of renormalization. Neither is a fine-tuning.

\paragraph{Distinction from $\Lambda$CDM.}
In $\Lambda$CDM, a cosmological constant $\Lambda$ is inserted by hand to account
for the observed accelerated expansion, with its value fine-tuned to $\rho_\Lambda
\sim 10^{-47}~\mathrm{GeV}^4$ against the QFT vacuum prediction of $\sim
10^{74}~\mathrm{GeV}^4$. RAGC eliminates the $\sim\!10^{120}$ discrepancy
structurally: the vacuum does not gravitate because the vacuum contains no
actualized states. The effective expansion of the universe is then driven entirely
by the intrinsic graph growth mechanism of Prediction~1, not by vacuum energy.

\paragraph{Observational target.}
If vacuum energy does not contribute to the metric source, the CMB power spectrum
should reflect a universe whose expansion history is governed entirely by matter
and radiation actualization densities, without a $\Lambda$ term. The angular power
spectrum, particularly at low multipoles where dark energy effects are most
pronounced, should deviate from $\Lambda$CDM predictions in a manner calculable
from the RA-modified Einstein equation
$G_{\mu\nu} = 8\pi G\,\langle\psi_{\mathrm{phys}}|\Tact|\psi_{\mathrm{phys}}\rangle$.
Quantifying this deviation is a concrete observational target for the program,
pending the Bianchi compatibility derivation (Section~\ref{sec:hardwall}).

\subsection{Prediction 3: Cosmological Constant Selection in Daughter Spacetimes}

When the actualization flux approaches the bandwidth limit $\Rmax$, a topological
severance nucleates a causally disconnected daughter spacetime. The Ledger Partition
Theorem (Theorem~\ref{thm:cut}) establishes that the Local Ledger Condition holds
independently on each side of the cut. The energy density inherited by the daughter
spacetime at nucleation is governed by the boundary flux at that moment, which is
determined by the local matter density and interaction rates in the parent
spacetime.

\paragraph{The prediction.}
RAGC predicts that the effective cosmological constant $\Lambda_{\mathrm{eff}}$ of
a daughter spacetime is not a fundamental free parameter but a derived quantity:
\begin{equation}
  \Lambda_{\mathrm{eff}} \;\sim\; \Phi_{\mathrm{boundary}}(C^*),
\end{equation}
where $\Phi_{\mathrm{boundary}}(C^*)$ is the total actualization flux crossing the
severance boundary $C^*$ at the moment of nucleation. Spacetimes nucleated from
high-flux regions inherit higher $\Lambda_{\mathrm{eff}}$; those from low-flux
regions inherit lower values.

\paragraph{Distinction from the anthropic landscape.}
In the string theory landscape, $\Lambda$ is selected from an enormous statistical
ensemble of vacua, with the observed value explained anthropically. In RAGC, the
selection mechanism is dynamical and local: $\Lambda_{\mathrm{eff}}$ is determined
by a specific, in-principle-measurable physical quantity (the boundary flux) at a
specific event (the nucleation). This is a sharper, more falsifiable account: it
ties the cosmological constant to the large-scale structure and matter distribution
of the parent spacetime at the nucleation epoch.

\paragraph{Observational imprint.}
If our universe is a daughter spacetime in this sense, the nucleation event should
leave an imprint in the large-scale structure of the early universe---specifically,
in correlations between the matter distribution at the nucleation epoch and the
effective cosmological constant inferred from late-time observations. Identifying
this imprint is a long-range observational target.

\paragraph{Predicted parent black-hole mass.}
Inverting the horizon-entropy relation with the observed $\Lambda_{\mathrm{obs}}$
gives a structural prediction for the parent black-hole mass at nucleation:
\begin{equation}
  M_{\mathrm{parent}} \;\approx\; 6.6 \times 10^6\,M_\odot,
  \label{eq:SMBH_parent}
\end{equation}
consistent with the mass of Sagittarius~A* ($4.1\times 10^6\,M_\odot$)
within the systematic uncertainty on $\Lambda$.  This is not a fit: it is
the unique mass implied by the Causal Severance mechanism applied to the
observed cosmological constant.

\paragraph{Three-in-one cosmological correlation.}
The Causal Severance mechanism couples $\Lambda$, $\Omega_{\mathrm{CDM}}$,
and $H_0$ through the boundary actualization flux:
\begin{equation}
  \Lambda \times \Omega_{\mathrm{CDM}} \times H_0^2
  \;\propto\; \Phi_{\mathrm{boundary}}^2.
  \label{eq:three_in_one}
\end{equation}
A correlated shift in all three parameters in the ratio this equation
predicts is a falsifiable RA signature: $\Lambda$CDM treats them as
independently tunable, but RA links them through a single dynamical quantity.

\paragraph{Connection to biological complexity.}
The causal graph mechanisms of RAGC extend directly to the emergence of biological
complexity. The Local Ledger Condition and Ledger Partition Theorem, proved here,
ground the conservation structure that biological assembly inherits. The companion
paper RACI~\citep{Sandeman2026c} establishes Condition~C1$'$: each JOIN operation
in Assembly Theory's molecular construction program corresponds to at least one
irreversible actualization event, operationally defined as an entropy-producing
on-shell energy transfer not reversed by subsequent unitary evolution. This
coarse-grained lower bound is supported by a variational argument and a
B3LYP/6-311+G* computational survey across the five primary bond types of
biological assembly (C-C, C-N, C-O, S-S, P-O), all confirmed exothermic.
The physics that generates the DAG metric is the same physics that makes
biological complexity possible. The further extension to agentic intelligence---
formalizing how biological systems develop generative models of their causal
environment and use them to steer the growth of the DAG---is developed in the
companion paper RAHC~\citep{Sandeman2026RAHC}.

% ─────────────────────────────────────────────────────────────────────────────

\subsection{Prediction 4: The Dark Matter Signature as an Actualization-Density
Effect}
\label{sec:pred4}

A particle species that undergoes no on-shell electromagnetic or strong-force
interactions contributes negligible actualization flux to the causal DAG and
therefore negligible metric sourcing. RAGC predicts that Weakly Interacting
Massive Particles (WIMPs) cannot account for galactic gravitational lensing, and
that the observed ``dark matter'' signatures have two distinct actualization-density
explanations requiring no new particle species.

\paragraph{The WIMP Prohibition.}
For a canonical 100~GeV WIMP with spin-independent cross-section
$\sigma_\chi \lesssim 10^{-44}$~cm$^2$ (consistent with current LZ and XENON1T
bounds), the actualization rate in the Milky Way ISM
($n_b \sim 1$~nucleon~cm$^{-3}$, $v \sim 200$~km~s$^{-1}$) is:
\begin{equation}
  \Gamma_\chi = n_b \cdot \sigma_\chi \cdot v
  \;\approx\; 2 \times 10^{-37}~\text{s}^{-1}~\text{per WIMP.}
  \label{eq:wimp_rate}
\end{equation}
The mean time between actualization events is $\tau = 1/\Gamma_\chi \approx
5 \times 10^{36}$~s, which is $\sim\!10^{18}$ times the age of the universe.
A WIMP traverses the causal DAG without writing vertices.

By contrast, a baryonic proton in the same environment undergoes electromagnetic
actualization at:
\begin{equation}
  \Gamma_p = n_b \cdot \sigma_{\mathrm{EM}} \cdot v_p
  \;\approx\; 10^{-14}~\text{s}^{-1},
  \label{eq:proton_rate}
\end{equation}
using a conservative atomic Coulomb cross-section $\sigma_{\mathrm{EM}} \sim
10^{-20}$~cm$^2$. The ratio of volumetric actualization flux is:
\begin{equation}
  \frac{\lambda_{\mathrm{baryon}}}{\lambda_{\mathrm{WIMP}}}
  \;\approx\; 10^{25}\text{--}10^{38},
  \label{eq:flux_ratio}
\end{equation}
where the range reflects different galactic environments (diffuse ISM to
stellar-density regions). In RA, the metric is sourced by actualization flux
$\lambda(x)$, not rest-mass density. A WIMP halo of five times the baryonic mass
contributes $\lesssim 10^{-25}$ of the metric sourcing of the visible baryonic
matter. This is a categorical prohibition, not a quantitative suppression: WIMPs
cannot account for observed galactic gravitational lensing in the RA framework.

\paragraph{Flat rotation curves, lensing, and the DJW mechanism.}
The same actualization-density field that produces the WIMP prohibition also
accounts for the dark-matter phenomenology, via two complementary RA-derived
mechanisms developed in RADM~\citep{Sandeman2026RADM}:
\begin{enumerate}[noitemsep]
  \item \textbf{Inner disk ($r < r_c$):} The 2D cylindrical Laplacian of
  $\ln\lambda$ for an exponential disk profile gives a surface density
  $\Sigma_\lambda = \xi/(4\pi G R r_d)$, producing a perfectly flat rotation
  curve with $v_\mathrm{flat}^2 = \xi/(2r_d)$ (RADM~\S3.1, fully derived).
  \item \textbf{Sparse halo ($r > r_c$, $\mu < 1$):} The Durhuus-Jonsson-Wheater
  theorem gives spectral dimension $d_s = 2$, so gravity is logarithmic:
  $\Phi_\mathrm{DJW}(r) = v_\mathrm{flat}^2\ln r$. The null geodesic
  deflection integral $\hat\alpha = -2\pi v_\mathrm{flat}^2/c^2$ = const,
  giving lensing convergence $\kappa(b) = v_\mathrm{flat}^2/(4G\Sigma_\mathrm{cr}b)$
  (RADM~\S4.1, RA-native derivation from DJW causal structure).
\end{enumerate}
The covariant form $\Box(\ln\lambda)$ is the open hard wall; the weak-field
results above are complete and RA-native. The Bullet Cluster offset follows
from the $A_\mathrm{RA}$(stellar)/$ A_\mathrm{RA}$(gas) ratio of $\sim 10^8$
(RADM~\S4.2).

\subsection{Prediction 5: CMB Acoustic Peak Structure and the Dark Matter Density}
\label{sec:pred5}

Two new results close major gaps in RAGC's cosmological programme.

\paragraph{Sinc suppression at acoustic peaks.}
For a baryon perturbation $\delta_b \propto \cos(k c_s t)$, the
actualization-density perturbation $\delta_A$ satisfies:
\begin{equation}
  \frac{\delta_A}{\delta_b} \;=\; \frac{\sin(k r_s)}{k r_s},
  \label{eq:sinc_pred5}
\end{equation}
where $r_s$ is the sound horizon.  At every CMB acoustic peak
($k r_s = n\pi$), $\sin(n\pi)/(n\pi) = 0$ exactly: the actualization-density
perturbation vanishes at the peaks independent of any free parameter.
The $\rho_A$ power spectrum is therefore automatically CDM-like in shape at
all acoustic scales, with no tuning required.  This is a structural prediction
distinguishable from $\Lambda$CDM by the absence of a CDM mass parameter.

\paragraph{BDG derivation of $\Omega_{\mathrm{CDM}}/\Omega_b$.}
The companion papers RASM and ra\_bdg\_couplings~\citep{Sandeman2026RASM,Sandeman2026BDG}
derive the dark matter/baryon density ratio from the BDG integers alone:
\begin{align}
  \frac{\Omega_{\mathrm{CDM}}}{\Omega_b}
  &= \frac{W_{\mathrm{other}}}{W_{\mathrm{baryon}}} \times \alpha_s(2m_p)
   = 17.32 \times 0.312 = 5.40,
  \label{eq:f0_ragc}
\end{align}
where $W_{\mathrm{baryon}} = 3/2$ exactly and $W_{\mathrm{other}}/W_{\mathrm{baryon}} = 17.32$
are proved from BDG combinatorics at $\mu=1$, and $\alpha_s(2m_p) = 0.312\pm
0.004$ from FLAG~2021 lattice QCD running.  This agrees with Planck~2018
$5.416$ to $0.3\%$ --- the first derivation of a cosmological density ratio
from pure causal-graph topology.

\paragraph{RA-native derivation of $Q_{\mathrm{eff}}$ and $m_p$ (RASM).}
The companion paper RASM~\citep{Sandeman2026RASM} provides an RA-native
identification of the effective scale.  The BDG filtration length $L=4$ for
quarks (Lean-verified, RATM) and the baryon weight $W_{\mathrm{baryon}}=3/2$
give $Q_{\mathrm{eff}} = L\times W_{\mathrm{baryon}}\times\Lambda_{\mathrm{QCD}}
= 6\Lambda_{\mathrm{QCD}} = 2m_p$ (exact by $m_p=3m_0$ and $m_0=m_p/3$).
Furthermore the BDG Regge formula (RASM~\S7) gives $m_p = \sqrt{c_4}\,\Lambda_{\mathrm{QCD}}$
to $0.08\%$, so the complete $f_0$ chain
$f_0 = 17.32\times\alpha_s(2\sqrt{c_4}\,\Lambda_{\mathrm{QCD}}) = 5.38$
is fully RA-native once $\Lambda_{\mathrm{QCD}}$ is derived from the
BDG string tension $\sigma = \Lambda_{\mathrm{QCD}}^2$ (Derivation~2 of RAGC).

\section{Mathematical Derivations}
\label{sec:math}

The following four derivations represent the core mathematical open tasks of
RAGC. They are \emph{derivations}---mathematical tasks required to connect the
framework's postulates to quantitative predictions---not empirical questions.
The predictions of Section~\ref{sec:predictions} stand on the basis of the
postulates and proved theorems regardless of whether these derivations are complete.
For each derivation, we provide a partial result or proof sketch that takes the
formalism as far as currently tractable, identifies the hard wall precisely, and
states what a collaborator would need to do to complete the result.

% ─── Problem 1 ───────────────────────────────────────────────────────────────
\subsection{Problem 1: The Generative Algebra of the Causal Graph}
\label{sec:algebra}

\paragraph{Goal.} Define the rigorous algorithmic update rules by which the
universe builds its DAG topology, incorporating the on-shell constraint.

\subsubsection{Foundation: Classical Sequential Growth Dynamics}

The Classical Sequential Growth (CSG) dynamics of Rideout and
Sorkin~\citep{RideoutSorkin2000} defines a growing sequence of finite causal sets
$\mathcal{C}_{0} \subset \mathcal{C}_{1} \subset \mathcal{C}_{2} \subset \cdots$
via Markov transitions. At each step, one new element $e$ is added with transition
probability
\begin{equation}
  t\bigl(\mathcal{C}_{n} \to \mathcal{C}_{n} + \{e\}\bigr)
  \;=\;
  \prod_{k} c_{k} \cdot \binom{n_{k}}{\text{coeff}},
  \label{eq:amplitude}
\end{equation}
where $c_{k}$ are coupling constants over precursor sets and $n_{k}$ is the count
of elements with that precursor structure. The key property is \emph{causal
invariance}: the transition probabilities commute for spacelike-separated
additions, so the final distribution is independent of the order in which such
additions are made---the discrete analogue of diffeomorphism invariance.

\subsubsection{The RA Extension: On-Shell Filter}

RAGC augments CSG dynamics with a physical constraint. Each potential new element
$e$ is assigned a type label $\tau(e) \in \{\text{on-shell},\,\text{off-shell}\}$,
determined by whether the underlying quantum transition has support on the
mass-shell $\Sigmam$ (see Section~\ref{sec:pact}). The RA-extended transition
probability is
\begin{equation}
  t_{\mathrm{RA}}\bigl(\mathcal{C}_{n} \to \mathcal{C}_{n} + \{e\}\bigr)
  \;=\;
  t_{\mathrm{RS}}\bigl(\mathcal{C}_{n} \to \mathcal{C}_{n} + \{e\}\bigr)
  \cdot \delta_{\tau(e),\,\text{on-shell}}.
\end{equation}
The on-shell filter is local (per-element), so causal invariance is inherited from
Rideout-Sorkin: if two spacelike-separated transitions commute in the unfiltered
theory, they commute in the filtered theory because the filter acts independently
on each. The adjacent possible at step $n$ is
\begin{equation}
  \AP(\mathcal{C}_{n})
  \;:=\;
  \bigl\{\, \mathcal{C}_{n} \cup \{v\} :
    \tau(v) = \text{on-shell},\;
    \exists\,\text{precursor antichain } A \subseteq \mathcal{C}_{n}
  \bigr\}.
\end{equation}

\subsubsection{Process Algebra Formulation}

Following Sulis~\citep{Sulis2020}, we express the RA-CSG dynamics as a
CCS-like process algebra. The causal set process $P$ is defined by
\begin{align*}
  P_{\text{causal set}} \;::=\;&
    0 & &\text{(no further on-shell transitions)}\\
    \mid\;& \tau_{\text{on-shell}}.\,P_{\text{causal set}}
            & &\text{(on-shell transition, then continue)}\\
    \mid\;& P + Q
            & &\text{(non-deterministic choice of next element).}
\end{align*}
Actualization corresponds to the communication of two processes (emitter and
absorber) along a shared on-shell channel. This formalism makes the bilateral
nature of actualization explicit: a vertex nucleates only when both participating
systems contribute an on-shell event.

\subsubsection{Hard Wall: The Generative Measure and Benincasa-Dowker Filter}

\textbf{Note: the primary hard wall (continuum limit) is now closed.}
The companion paper RACL~\citep{Sandeman2026RACL} proves that the
Lean-verified Local Ledger Condition, combined with the Benincasa-Dowker
theorem, Lovelock uniqueness, and the $\mathcal{P}_{\mathrm{act}}$
conservation theorem, uniquely determines
$G_{\mu\nu} = 8\pi G\,\mathcal{P}_{\mathrm{act}}[T_{\mu\nu}]$ with
$\Lambda = 0$ and no free parameters. The BDG locality lemma in RACL
proves $\tau_{\mathrm{BDG}}(\lambda) = 4\Gamma(5/4)(12/\pi^2)^{1/4}
\lambda^{-1/4} \to 0$ as $\lambda \to \infty$, confirming that
the dense-graph limit is local and second-order.
The remaining open problem is fixing the coupling constants $c_k$ in the RA-CSG
transition probabilities. In standard Rideout-Sorkin CSG, these are genuinely
free parameters~\citep{Sorkin2003}. RA provides a principled discrete constraint
that vanilla CST lacks: the actualization filter.

\paragraph{The Benincasa-Dowker sign criterion.}
The Benincasa-Dowker-Glaser (BDG) action~\citep{Benincasa2010} provides the
unique discrete, retarded, Lorentz-invariant action functional for a causal set.
At a candidate vertex $v_*$ in 4D, it is pure graph combinatorics:
\begin{equation}
  S_{\mathrm{BDG}}^{(4)}(v_*)
  = \frac{1}{l_P^2}
    \bigl(1 - N_1(v_*) + 9N_2(v_*) - 16N_3(v_*) + 8N_4(v_*)\bigr),
  \label{eq:BDG}
\end{equation}
where $N_k(v_*)$ counts the $k$-element inclusive causal intervals in the
immediate past of $v_*$. No metrics or PDEs appear; only the combinatorial
structure of the causal past is used. In the continuum limit, the mean of
$S_{\mathrm{BDG}}$ recovers the Einstein-Hilbert action~\citep{Benincasa2010}.

The RA actualization criterion $\Delta S(\rho\|\sigma_0) > 0$ has a natural
discrete counterpart via the BDG action: a candidate vertex $v_*$ actualizes
if and only if
\begin{equation}
  S_{\mathrm{BDG}}^{(4)}(v_*) > 0.
  \label{eq:BDG_filter}
\end{equation}
This identification is RA-native: a vertex contributes positively to the local
scalar curvature if and only if it increases the local relative entropy with
respect to the vacuum. The sign criterion replaces the phenomenological
on-shell filter with a purely graph-theoretic one.

\paragraph{Mean-field approximation and Poisson-CSG.}
The RA actualization filter, combined with QFT cluster decomposition
(spacelike-separated interactions decohere independently), implies that the
number of actual causal connections $k$ formed by a new vertex in an environment
with actualization density $\lambda(x)$ follows
\begin{equation}
  c_k^{\mathrm{RA}}(x) = \frac{e^{-\mu(x)}\,\mu(x)^k}{k!},
  \qquad
  \mu(x) = \lambda(x)\,V_{\mathrm{coh}}(x)\,p_{\mathrm{th}},
  \label{eq:poisson_ck}
\end{equation}
where $V_{\mathrm{coh}}(x)$ is the coherence volume of the local quantum state
and $p_{\mathrm{th}} = \alpha_{\mathrm{EM}} \approx 1/137 \approx 0.0073$ is the Kinematic Coherence Bound of
RAQI~\citep{Sandeman2026RAQI}. This Poisson-CSG formula is the mean-field
approximation to the exact BDG filter. It makes $c_k$ \emph{local and
environment-dependent} --- a structural difference from all prior CSG
formulations, where $c_k$ are universal constants.

This locality has two immediate consequences. In dense environments
($\mu \gg 1$): $c_k$ peaks at large $k$, the graph is highly connected, $d_s
= 3$, and standard GR is recovered. In sparse environments ($\mu \ll 1$): $c_k$
is dominated by $k = 0$ and $k = 1$, the graph becomes tree-like, and by the
Durhuus-Jonsson-Wheater theorem~\citep{Durhuus2010}, $d_s = 2$. This is the
mechanism that drives dimensional reduction in galactic halos
(RADM~\citep{Sandeman2026RADM}, Section~3.2) and underlies the conditional
Tully-Fisher derivation.

\paragraph{Remaining gap.}
Two quantities require determination before the Poisson-CSG formula yields
quantitative predictions:
\begin{enumerate}
\item The coherence volume $V_{\mathrm{coh}}(x)$: the volume of the
  pre-actualization quantum state (the off-shell entanglement region that
  collapses upon a Kinematic Snap). This is not the Planck volume but depends
  on the local energy scale and quantum state structure.
\item The exact $c_k$ distribution beyond the mean-field limit: the full BDG
  + RA filter calculation requires MCMC simulation of the RA-CSG growing under
  equation~\eqref{eq:BDG_filter}, extracting the asymptotic $c_k$ distribution
  and the ratio $\bar\lambda_m/\bar\lambda_{\mathrm{exp}}$ (which determines
  the Hubble tension prediction of Section~\ref{sec:predictions}).
\end{enumerate}
\textbf{Collaborator target:} causal set theory, CDT computational physics,
quantum information theory.

% ─── Problem 2 ───────────────────────────────────────────────────────────────
\subsection{Problem 2: The DAG Continuum Limit as an Emergent Shadow}
\label{sec:continuum}

\paragraph{Goal.} Prove that the continuous equations of General Relativity are
the emergent shadows of the discrete graph-rewriting process in the limit of
infinite vertex density.

\subsubsection{LLC \texorpdfstring{$\to$}{→} Continuity Equation}

Working in natural units ($c = \hbar = 1$) throughout this section.
Let $G$ be a DAG with vertex density $\lambda$ (vertices per unit 4-volume, units
$[\text{length}]^{-4}$) satisfying the LLC. Define the smoothed charge current
\begin{equation}
  J^{\mu}_{\mathrm{smooth}}(x)
  \;:=\;
  \frac{1}{\varepsilon^{4}\,\lambda}
  \sum_{\substack{v:\,|v-x|<\varepsilon}} q(v)\,u^{\mu}(v),
\end{equation}
where $u^{\mu}(v)$ is the unit future-directed tangent to the local causal
direction at $v$. The denominator $\varepsilon^{4}\lambda$ counts the expected
number of vertices in the averaging 4-volume---it is dimensionless---so
$J^{\mu}_{\mathrm{smooth}}$ has units of charge density current in natural units
($[\text{charge}]\cdot[\text{length}]^{-3}$), consistent with the standard
4-current. As $\lambda \to \infty$ and $\varepsilon \to 0$ with
$\varepsilon^{4}\lambda \to \mathrm{const}$, a standard Taylor expansion on the
discrete divergence shows
\begin{equation}
  \partial_{\mu}J^{\mu}_{\mathrm{smooth}} \;\to\; 0 + \mathcal{O}(1/\lambda),
\end{equation}
recovering the continuity equation $\nabla_{\mu}J^{\mu} = 0$. The error term
$\mathcal{O}(1/\lambda)$ vanishes in the continuum limit. This argument is
structurally identical to the Bombelli-Henson-Sorkin recovery of the d'Alembertian
from a causal set \citep{Bombelli2009}, adapted to the current vector-field context.

\subsubsection{Metric Reconstruction}

The Hawking-Malament theorem establishes that on a Lorentzian manifold, the causal
structure plus a volume measure uniquely determines the metric (up to conformal
factor). In the DAG model: the causal structure is the graph $G$ itself, and the
volume measure is the vertex count. Therefore, in principle, the metric is
determined by the DAG once the continuum limit is taken.

Whether the RA-CSG dynamics generates a DAG that, in the continuum limit, satisfies
$G_{\mu\nu} = 8\pi G\,\langle\Tact\rangle$ rather than some other metric is the
core open question. Causal Dynamical Triangulation (CDT) \citep{Ambjorn2012}
provides a benchmark: CDT recovers 4D de Sitter spacetime from a purely
combinatorial path integral, establishing that the continuum recovery program is
achievable from discrete causal structures. The analogous result for RA-CSG is
\emph{conjectured but not proven}.

\subsubsection{The Unruh Reconciliation: Precise Problem Statement and
Resolution Strategy}

The Unruh effect presents a conceptual tension for RA. An accelerated (Rindler)
observer detects a thermal bath of particles in the Minkowski vacuum. In RA terms:
the inertial observer has $\Pact\ket{0}_{\mathrm{Mink}} = 0$ (no on-shell
excitations). The Rindler observer sees a thermal state with $\langle
T_{\mathrm{act}}\rangle \neq 0$ relative to their vacuum. Does this make
actualization observer-dependent?

\paragraph{Why particle-number cannot be the fundamental criterion.}
The red-team formulation of this objection is precise: if $\Pact$ is defined in
terms of Fock-space particle content, it is frame-dependent, because Rindler and
Minkowski observers disagree on particle number by definition. This is correct,
and it means that the Fock-space definition of $\Pact$ in
Section~\ref{sec:pact} must be understood as the weakly-coupled, perturbative
limiting expression of a more fundamental criterion---not the criterion itself.

\paragraph{The AQFT resolution: relative entropy as the fundamental criterion.}
The fundamental actualization criterion in RA is not ``a particle is emitted''
but rather ``an irreversible increase in entanglement entropy between a
localised system and its asymptotic environment occurs.'' In the language of
algebraic QFT, this is a statement about \emph{relative entropy with respect to
the vacuum state}---a quantity that is diffeomorphism-invariant and
observer-independent.

Formally: an actualization event occurs when an interaction induced by the
local algebra $\mathcal{A}(\mathcal{O})$ produces a strictly irreversible
increase in the relative entropy $S(\rho \| \sigma_0)$, where $\rho$ is the
post-interaction state and $\sigma_0$ is the vacuum state. Relative entropy is a
scalar under diffeomorphisms. A strictly irreversible thermodynamic exchange is
an objective, coordinate-independent fact.

The Bisognano-Wichmann theorem~\citep{BisognanoWichmann1975} in algebraic QFT
establishes that the Rindler and Minkowski algebras are unitarily
equivalent---describing different perspectives on the same physics, not different
physics. The Minkowski and Rindler observers may disagree on the \emph{particle
content} of an interaction, but they must agree on whether a thermodynamically
irreversible event occurred, because irreversibility is encoded in the
monotone decrease of relative entropy---a frame-independent quantity. That
irreversible event registers as a permanent vertex on the DAG regardless of which
observer describes it.

\paragraph{Finite-dimensional proof and remaining AQFT target.}
The core claim of frame-independence has been proved in finite dimensions and
formally verified in Lean~4. Specifically, the theorem
\texttt{frame\_independence} in \texttt{RA\_AQFT\_Proofs\_v10.lean} establishes
that $S(U\rho U^\dagger \| \sigma_0) = S(\rho \| \sigma_0)$ for any unitary $U$
representing a Lorentz boost, where the proof uses unitary invariance of relative
entropy together with the Poincar\'{e}-invariance of the vacuum as a physics
axiom. The key lemma \texttt{log\_unitary\_conj}---that $\log(UMU^\dagger) =
U\log(M)U^\dagger$ for Hermitian $M$---is a one-line corollary of
\texttt{Matrix.cfc\_conj\_unitary}, the general CFC unitary-conjugation result
from the Lean-QuantumInfo library~\citep{Meiburg2025LQI}.

The remaining target at the full AQFT level is the infinite-dimensional
extension: mapping the finite-dimensional relative entropy proof to Araki's modular
relative entropy~\citep{Araki1976} for type~III$_1$ von Neumann algebras, where
standard density matrices and the trace definition of relative entropy do not exist.
The precise target is to prove that the Rindler KMS state exhibits
$\Delta S_{\mathrm{rel}} = 0$ in the sense of Araki's relative entropy on the
wedge algebra $\mathcal{A}(W)$, and that $\Pact$ defined on the local net commutes
with the Tomita-Takesaki modular flow $\sigma_s^{\Omega}$ of the vacuum state
$\Omega$. This is the exact gap between the finite-dimensional Lean result and a
rigorous QFT theorem. \textbf{A collaborator with expertise in Tomita-Takesaki
modular theory and Araki's relative entropy would be well-positioned to complete
this extension.}

\paragraph{The stationarity criterion: why the Rindler thermal state does not actualize.}
The Kinematic Snap provides a precise answer to the Unruh question that does not
require the full AQFT proof. A uniformly accelerated observer sees the Minkowski
vacuum as a thermal bath at temperature $T_U = \hbar a / (2\pi c k_B)$. In the
correct algebraic QFT language, the relevant object is the \emph{restriction of
the Minkowski vacuum to the Rindler wedge algebra}---a KMS state with respect to
boost automorphisms (the Bisognano-Wichmann theorem~\citep{BisognanoWichmann1975}).
This KMS (thermal) state is in equilibrium with respect to modular flow: its
relative entropy with respect to the vacuum reference state is
\emph{constant in time} ($\Delta S_{\mathrm{rel}} = 0$).
No irreversible increase in relative entropy is occurring. No net on-shell bosons
are being emitted into the asymptotic environment at a rate exceeding absorption.
The Rindler KMS state is the universe's stationary potentia as seen by the
accelerated observer---a persistent configuration of the quantum field, not a
sequence of actualization events.

Both inertial and accelerated observers therefore agree on the causal fact: in the
Minkowski vacuum, no actualization is occurring. They disagree only on the modal
decomposition (which modes to call ``particles'')---a difference of description,
not a difference of physical reality. Actualization requires an irreversible
\emph{increase} in relative entropy, not merely non-zero relative entropy.

\paragraph{The non-unitary element.}
A natural question arises: since unitary evolution preserves relative entropy,
how can $\Delta S_{\mathrm{rel}} > 0$ ever occur? The answer is that actualization
is precisely the non-unitary element. The Kinematic Snap is an open-system event:
the on-shell boson propagates to the asymptotic environment, which is traced over
in forming the reduced state of the system. This environmental tracing---the
coarse-graining implicit in the Local Ledger Condition---is the non-unitary
step that permits irreversible entropy increase. In the language of open quantum
systems, the actualization map is a CPTP (completely positive trace-preserving)
channel, not a unitary: it is precisely the $\mathcal{E}$ in the data processing
inequality $S(\mathcal{E}(\rho) \| \mathcal{E}(\sigma)) \leq S(\rho \| \sigma)$.
An actualization event is the physical realisation of such a channel, with the
environment providing the non-unitary element. Formalising $\Pact$ as a specific
CPTP map on the local net $\mathcal{A}(\mathcal{O})$ and proving its modular
covariance is the open target identified in D2.

\paragraph{Note on virtual exchanges and relative entropy.}
A claim appearing in the supporting computational scripts---that ``off-shell
(virtual) exchanges have $S(\rho_{\mathrm{off-shell}} \| \sigma_0) = 0$''---
requires clarification. Virtual exchanges are internal-line bookkeeping
in perturbation theory; they are not standalone density operators on the
wedge algebra in the QFT sense. The correct statement is operational: interactions
that remain kinematically constrained to the virtual (off-shell) regime produce
no asymptotic on-shell bosons, no environmental record, and therefore no
irreversible increase in $S(\rho \| \sigma_0)$. The absence of a DAG vertex
is not equivalent to the system being in a zero-relative-entropy state; it means
no irreversible record-forming event has occurred. The relative entropy criterion
is meaningful for post-interaction states where a comparison to the vacuum can be
made; the ``virtual regime'' is precisely the regime in which no such comparison
is yet applicable.

The finite-dimensional Lean proof of \texttt{rindler\_stationarity} confirms this
picture formally: the theorem establishes that $S(U\rho_R U^\dagger \| \sigma_0) =
S(\rho_R \| \sigma_0)$ for any Lorentz boost $U$ applied to the Rindler thermal
state, which is precisely the statement that $\Delta S_{\mathrm{rel}} = 0$.
The remaining step at the full AQFT level is to extend this result to the
infinite-dimensional type~III$_1$ algebra setting, replacing the finite-dimensional
relative entropy with Araki's modular relative entropy and proving that the Rindler
KMS state is a fixed point of the modular flow in that setting.

% ─── Problem 3 ───────────────────────────────────────────────────────────────
\subsection{Problem 3: \texorpdfstring{$\Pact$}{P\_act} as a Syntactic Type-Constraint}
\label{sec:pact}

\paragraph{Goal.} Define the Actualization Projector rigorously, prove that it
suppresses the vacuum energy contribution to the metric source term, and establish
its gauge covariance.

\subsubsection{The Standard Vacuum Energy Problem}

In standard QFT on a curved background, the Einstein equation takes the form
\begin{equation}
  G_{\mu\nu} + \Lambda g_{\mu\nu}
  \;=\;
  8\pi G\,\langle 0 | T_{\mu\nu} | 0 \rangle.
\end{equation}
The QFT vacuum contribution $\langle 0|T_{\mu\nu}|0\rangle \approx M_{\mathrm{Pl}}^{4}
\sim 10^{74}~\mathrm{GeV}^{4}$, while observations require $\rho_{\Lambda} \sim
10^{-47}~\mathrm{GeV}^{4}$---a discrepancy of $\sim\!10^{120}$. The standard approach
renormalizes by subtracting the vacuum term as a counterterm; this is technically
consistent but requires extraordinary fine-tuning \citep{Weinberg1989}.

\subsubsection{Definition of \texorpdfstring{$\Pact$}{P\_act}}

Let $\mathcal{H}$ be the Hilbert space of the quantum field. Define the on-shell
subspace
\begin{equation}
  \Hon
  \;:=\;
  \bigl\{\,\ket{\psi} \in \mathcal{H} :
    p^{\mu}p_{\mu}\ket{\psi} = m^{2}c^{2}\ket{\psi}
  \bigr\}.
\end{equation}
In momentum space, $\Hon$ is supported on the positive-frequency mass hyperboloid
$\Sigmam := \{k^{\mu} \in \mathbb{R}^{3,1} : k^{\mu}k_{\mu} = m^{2}c^{2},\;
k^{0} > 0\}$.

\begin{definition}[Actualization Projector]
$\Pact$ is the spectral projection onto $\Hon$ (natural units $c = \hbar = 1$
throughout):
\begin{equation}
  \Pact
  \;:=\;
  \int_{\Sigmam} \frac{d^{3}k}{(2\pi)^{3}}\,\frac{1}{2\omega_{\mathbf{k}}}
  \;\ket{k}\!\bra{k},
\end{equation}
where $\omega_{\mathbf{k}} = \sqrt{|\mathbf{k}|^{2} + m^{2}}$ is the
on-shell energy (in natural units the dispersion relation is
$p^{\mu}p_{\mu} = m^{2}$, so $\omega_{\mathbf{k}}^{2} = |\mathbf{k}|^2 + m^2$),
and the integration is over the Lorentz-invariant measure on $\Sigmam$.
The on-shell subspace definition uses $p^{\mu}p_{\mu} = m^{2}$ throughout;
in SI units this reads $p^{\mu}p_{\mu} = m^{2}c^{2}$ and
$\omega_{\mathbf{k}} = \sqrt{|\mathbf{k}|^{2}c^{2} + m^{2}c^{4}}/\hbar$.
For a multi-species theory:
$\Pact = \sum_{i} \int_{\Sigma_{m_{i}}} \frac{d^{3}k}{(2\pi)^{3}}\,
\frac{1}{2\omega_{\mathbf{k},i}} \ket{k,i}\!\bra{k,i}$.
\end{definition}

\begin{proposition}
$\Pact$ is an orthogonal projection: $\Pact^{2} = \Pact$ and
$\Pact^{\dagger} = \Pact$.
\end{proposition}
\begin{proof}
Each $\ket{k}$ is an eigenstate of the mass operator $p^{\mu}p_{\mu}$ with
eigenvalue $m^{2}c^{2}$. The projection onto an eigenspace of a self-adjoint
operator is idempotent and self-adjoint by the spectral theorem.
\end{proof}

\subsubsection{Vacuum Energy Suppression}

\begin{definition}
The RA-modified stress-energy operator is $\Tact := \Pact\,T^{\mu\nu}\,\Pact$
(the compression of $T^{\mu\nu}$ to the on-shell subspace: $\bra{\psi}\Tact\ket{\phi} = 0$
unless both $\ket{\psi}$ and $\ket{\phi}$ have support on $\Sigmam$).
The modified Einstein equation is
\begin{equation}
  G_{\mu\nu}
  \;=\;
  8\pi G\,\langle\psi_{\mathrm{phys}}|\,\Tact\,|\psi_{\mathrm{phys}}\rangle,
\end{equation}
with no bare $\Lambda$ term.
\end{definition}

\begin{theorem}[Vacuum Energy Suppression]
\label{thm:vacuum}
$\langle 0|\,\Tact\,|0\rangle = 0$.
\end{theorem}
\begin{proof}
The QFT vacuum $\ket{0}$ is the Fock state with no on-shell particles. In the Fock
representation, $\Pact\ket{0} = 0$: $\ket{0}$ has no support on $\Sigmam$ because
it contains no on-shell quanta. Therefore
\begin{equation*}
  \langle 0|\,\Tact\,|0\rangle
  \;=\;
  \langle 0|\,\Pact\,T^{\mu\nu}\,\Pact\,|0\rangle
  \;=\;
  \langle 0|\,\Pact\,T^{\mu\nu}\cdot 0
  \;=\; 0. \qedhere
\end{equation*}
\end{proof}

At the ontological level, this result is a structural consequence of the RA
postulate: the vacuum is a superposition of unactualized potentia, and $\Pact$
removes it from the metric source by construction. It is not a fine-tuning.
At the operational level in semiclassical gravity, the same result is implemented
by the vacuum-subtraction prescription (Section~\ref{sec:pact}, hard wall~iv):
\begin{equation*}
  \langle\Tact\rangle[g,\Phi]
  := \langle T^{\mu\nu}_{\mathrm{ren}}\rangle[g,\Phi]
   - \langle T^{\mu\nu}_{\mathrm{ren}}\rangle[g,\ket{0}].
\end{equation*}
These are not competing claims: the ontological framing explains \emph{why} the
vacuum does not gravitate in RA; the prescription is the \emph{operational
implementation} in the renormalized QFT setting, ensuring that configuration
differences (Casimir stresses, binding energies) are correctly preserved.
The prescription is locally covariant in the flat-space regime; its extension
to curved backgrounds via Hadamard renormalization is an open task (D3 hard wall iv).
This is distinct from, and complementary to, the Kaloper-Padilla sequestering
mechanism~\citep{KaloperPadilla2014}; a future derivation should examine whether
they agree at leading order.

\subsubsection{Gauge Covariance}

\begin{proposition}
The on-shell condition $k_{\mu}k^{\mu} = g_{\mu\nu}k^{\mu}k^{\nu} = m^{2}c^{2}$
is a scalar under diffeomorphisms. Therefore $\Pact$ is well-defined in curved
spacetime at each fiber of the cotangent bundle.
\end{proposition}
\begin{proof}
Under a diffeomorphism $\varphi: M \to M$, both $k^{\mu}$ and $g_{\mu\nu}$
transform as tensors, so their contraction is preserved. The subspace $\Hon$ is
therefore defined by a coordinate-invariant condition.
\end{proof}

\subsubsection{Connection to the DAG Type-Constraint}

\begin{definition}[Vertex Nucleation Condition]
A quantum transition $T: v_{\mathrm{in}} \to v_{\mathrm{out}}$ nucleates a new
vertex $v \in V(G)$ if and only if $\exists\,k^{\mu}$ such that $k^{\mu}k_{\mu} =
m^{2}c^{2}$---equivalently, if $\Pact\ket{T} \neq 0$.
\end{definition}

\begin{claim}[Virtual Transitions Do Not Nucleate]
Off-shell internal lines in Feynman diagrams satisfy $\Pact\ket{T_{\mathrm{virt}}}
= 0$ and therefore contribute no vertices to the causal DAG.
\end{claim}

\begin{proof}[Supporting argument]
Internal Feynman lines have propagator $\Delta_{F}(k) = i/(k^{2} - m^{2} +
i\varepsilon)$ with $k^{2} \neq m^{2}$ generically. They represent superpositions
over the full off-shell spectrum, with no support on $\Sigmam$. Only asymptotic
(external) states and on-shell resonances (at the propagator pole $k^{2} = m^{2}$)
have support in $\Hon$. This is precisely the content of the LSZ reduction formula:
physical $S$-matrix amplitudes are the residues at on-shell poles. The causal DAG
is therefore built from $S$-matrix poles, not from the full off-shell generating
functional.
\end{proof}

\subsubsection{Hard Wall}
\label{sec:hardwall}

Three open sub-problems were identified in prior development of this framework.
A fourth, raised by external critique, is added here and is arguably the most
fundamental.

\begin{enumerate}[label=(\roman*)]
  \item \textbf{Non-perturbative regime.} The definition in
  Section~\ref{sec:pact} uses Fock-space structure, which requires a perturbative
  vacuum. In strongly coupled regimes (QCD confinement, condensed matter critical
  points), the quasi-particle picture breaks down. A non-perturbative
  definition---perhaps via algebraic QFT (Haag-Kastler axioms) or the OPE---is
  required.

  \item \textbf{Renormalization compatibility.} $\Tact$ may require its own
  renormalization scheme. The standard renormalization of $T^{\mu\nu}$ subtracts
  the vacuum contribution as a counterterm; $\Pact$ achieves this structurally.
  The two procedures must be shown to agree at the level of renormalized physical
  amplitudes. Moreover, even granting that $\langle 0|\Tact|0\rangle = 0$
  eliminates the absolute vacuum energy density from the metric source, a complete
  solution to the cosmological constant problem requires showing how
  \emph{changes} in vacuum energy---Casimir stresses, binding energy
  contributions, loop corrections to the effective potential---are treated by the
  RA-modified source. These are finite, physical, and measurable contributions
  that gravitate in standard semiclassical gravity; demonstrating that $\Pact$
  handles them consistently, or identifying the prediction that deviates from
  standard GR, is an explicit open task.

  \paragraph{The trace anomaly.}
  A related and sharper question concerns the conformal (trace) anomaly: the
  non-vanishing of $\langle T^{\mu}_{\ \mu}\rangle$ for classically
  conformally-invariant fields on curved backgrounds, which is essential for
  deriving Hawking radiation. The RA prescription preserves the trace anomaly
  because $\Pact$ acts on the \emph{already-renormalized} stress tensor
  $\langle T^{\mu\nu}_{\mathrm{ren}}\rangle$: the trace anomaly arises from the
  UV renormalization procedure (point-splitting or dimensional regularization)
  before the on-shell projection is applied. $\Pact$ then projects onto the
  on-shell sector of the renormalized theory, inheriting the trace anomaly from
  that theory. Consequently, the RA prescription does \emph{not} eliminate virtual
  loop contributions that generate the trace anomaly---it eliminates only the
  configuration-independent absolute vacuum energy. This is consistent with
  Hawking radiation: the thermal emission from a black hole is sourced by the
  trace anomaly (and related vacuum polarization effects) of the renormalized
  theory, which $\Pact$ preserves. The cosmological constant suppression and
  Hawking radiation are therefore not in tension within the RA framework.

  \item \textbf{Back-reaction self-consistency.} The modified metric $g_{\mu\nu}$
  (sourced by $\Tact$) feeds back into the definition of $\Hon$ via
  $g_{\mu\nu}k^{\mu}k^{\nu} = m^{2}c^{2}$. Establishing self-consistency
  requires a fixed-point argument.

  \item \textbf{Bianchi compatibility and local conservation.}
  This was the most fundamental open problem. It has been substantially
  resolved by identifying the correct RA prescription for the
  stress-energy source, as detailed below.

  \paragraph{The vacuum-subtraction prescription and Casimir benchmark.}
  The resolution rests on a precise distinction. The RA actualization
  projector $\Pact$ suppresses the \emph{absolute} vacuum energy---the
  configuration-independent contribution $\sim M_{\mathrm{Pl}}^4$ that
  is present in all states equally. It does \emph{not} suppress energy
  \emph{differences} arising from physically actualized boundary
  conditions, because those differences reflect real causal facts about
  the system. The Casimir effect is precisely such a difference:
  $E_{\mathrm{Casimir}} = E(\text{plates}) - E(\text{no plates})$. The
  plates are actualized objects that genuinely modify the causal graph;
  their boundary conditions are physically real.

  This gives the RA prescription a precise mathematical form:
  \begin{equation}
    \langle\Tact\rangle[g,\Phi]
    \;:=\;
    \langle T^{\mu\nu}_{\mathrm{ren}}\rangle[g,\Phi]
    \;-\;
    \langle T^{\mu\nu}_{\mathrm{ren}}\rangle[g,\ket{0}],
    \label{eq:ra_prescription}
  \end{equation}
  where $g$ is the spacetime metric, $\Phi$ are the matter fields, and
  $\ket{0}$ is the Minkowski (or de Sitter) vacuum state.

  \begin{theorem}[Bianchi Compatibility --- Casimir Geometry]
  \label{thm:bianchi}
  Under the vacuum-subtraction prescription~(\ref{eq:ra_prescription}):
  \begin{enumerate}[label=\emph{(\alph*)}]
    \item $\langle 0|\Tact|0\rangle = 0$ by construction.
    \textup{[Vacuum suppression]}
    \item $\nabla_{\mu}\langle\Tact\rangle = \nabla_{\mu}\langle
    T^{\mu\nu}_{\mathrm{ren}}\rangle - \nabla_{\mu}\langle
    T^{\mu\nu}_{\mathrm{ren}}\rangle\big|_{\ket{0}} = 0$, provided the
    standard renormalized tensor satisfies conservation under the Wald
    axioms (point-splitting scheme). \textup{[Bianchi compatibility]}
    \item The Casimir stress-energy is preserved exactly:
    $\langle\Tact\rangle(\text{plates}) = \langle
    T^{\mu\nu}_{\mathrm{Casimir}}\rangle$. \textup{[Physical vacuum
    effects reproduced]}
    \item The modified Einstein equation $G_{\mu\nu} = 8\pi G\,\langle
    \Tact\rangle$ is consistent: $\nabla_{\mu}G^{\mu\nu} \equiv 0$ and
    $\nabla_{\mu}\langle\Tact\rangle = 0$ both hold.
    \textup{[Field equation consistency]}
  \end{enumerate}
  \end{theorem}
  \begin{proof}
  (a)~Self-subtraction. (b)~Linearity of $\nabla_\mu$ and conservation
  of each term separately under Wald axioms. (c)~By definition:
  $\langle T^{\mu\nu}_{\mathrm{ren}}\rangle(\text{plates}) -
  \langle T^{\mu\nu}_{\mathrm{ren}}\rangle(\text{Minkowski}) =
  \langle T^{\mu\nu}_{\mathrm{Casimir}}\rangle$. (d)~Follows from
  (b) and the contracted Bianchi identity $\nabla_\mu G^{\mu\nu}
  \equiv 0$.
  \end{proof}

  \paragraph{Scope of the theorem.}
  Theorem~\ref{thm:bianchi} establishes Bianchi compatibility within the
  flat-spacetime regime. Step~(b) relies on the linearity of $\nabla_\mu$ and the
  assumption that each term separately satisfies $\nabla_\mu \langle
  T^{\mu\nu}_{\mathrm{ren}}\rangle = 0$ under the Wald axioms (point-splitting
  scheme). This is established for free quantum fields on globally hyperbolic
  spacetimes~\citep{Wald1994}. On general curved backgrounds, the prescription
  requires a curved-spacetime renormalization scheme (Hadamard point-splitting
  rather than heuristic vacuum subtraction), and the numerical benchmark uses
  ordinary derivatives $\partial_\mu$ in flat space rather than the covariant
  divergence $\nabla_\mu$ in a non-trivial metric. Two things follow from this:

  \begin{enumerate}[label=(\roman*),noitemsep]
    \item The flat-space Casimir benchmark is valid and numerically confirms the
    prescription reproduces all measurable vacuum physics correctly.
    \item Demonstrating $\nabla_\mu\langle\Tact\rangle = 0$ on a curved
    background---for example, the Casimir apparatus in a weak gravitational
    field---requires a full curved-spacetime Hadamard renormalization and constitutes
    a decisive test case (see D3 checklist in Section~\ref{sec:hardwall}).
  \end{enumerate}

  A numerical benchmark at $d = 1\,\mu$m gives Casimir pressure
  $\langle T^{zz}_{\mathrm{RA}}\rangle = -1.30 \times 10^{-3}$~Pa,
  consistent with the experimental value $\approx -1.3 \times
  10^{-3}$~Pa. The prescription reproduces all measurable vacuum effects
  while eliminating the $\sim\!10^{120}$ vacuum energy catastrophe.

  \paragraph{The emergent metric defence.}
  A deeper resolution follows from the ontological hierarchy of the RA
  framework. The Einstein field equations are not fundamental in RA;
  they are the thermodynamic equation of state for the underlying causal
  graph, in the same sense as Jacobson's thermodynamic derivation of
  GR~\citep{Jacobson1995}. The Local Ledger Condition enforces
  energy-momentum conservation exactly at every discrete actualization
  vertex. The continuous Bianchi identity $\nabla_\mu T^{\mu\nu} = 0$
  is the coarse-grained continuum limit of the LLC, not an independent
  constraint that the discrete theory must satisfy separately. In this
  sense, the Bianchi compatibility problem dissolves at the level of the
  fundamental theory: conservation is built into the graph architecture,
  and the continuous field equations inherit it in the continuum limit.

  \paragraph{Remaining open scope.}
  The vacuum-subtraction prescription assumes the standard
  $\langle T^{\mu\nu}_{\mathrm{ren}}\rangle[g]$ satisfies the Wald
  axioms, which hold for free fields and established interacting
  theories. The Wald ambiguity (freedom to add local conserved geometric
  counterterms) is fixed uniquely in the flat-space limit by the RA
  condition $\langle 0|\Tact|0\rangle = 0$. Extending to the full
  non-perturbative regime requires the algebraic QFT formulation of
  Derivation~2 (D2). These remaining gaps are the same gaps that exist
  in standard semiclassical gravity; RA does not introduce new
  difficulties beyond those already present in QFT on curved spacetime.\end{enumerate}

% ─── Problem 4 ───────────────────────────────────────────────────────────────
\subsection{Problem 4: The Graph Cut Theorem and Ledger Partitioning}

\paragraph{Goal.} Prove that the LLC holds independently on each side of a causal
severance, and characterize the implications for event horizons.

\subsubsection{Formal Setup}

Let $G = (V, E)$ be a DAG satisfying the LLC at every vertex. Let $C \subset E$ be
a \emph{cut}---a set of edges whose removal disconnects $G$ into two acyclic
subgraphs $G_{L} = (V_{L}, E_{L})$ and $G_{R} = (V_{R}, E_{R})$ with $V_{L} \cap
V_{R} = \emptyset$, $V_{L} \cup V_{R} = V$.

By the acyclicity of $G$ (global causal order), edges in $C$ run from $V_{L}$ to
$V_{R}$ (past to future). Boundary vertices in $V_{L}$ have outgoing cut-edges
only; boundary vertices in $V_{R}$ have incoming cut-edges only.

\subsubsection{Theorem and Proof}

\begin{theorem}[Ledger Partition]
\label{thm:cut}
Under the above setup:
\begin{enumerate}[label=\emph{(\alph*)}]
  \item Each subgraph $G_{L}$ and $G_{R}$ satisfies the LLC at all interior
  vertices (those with no edges in $C$).
  \item At boundary vertices, the LLC holds with cut-edges accounted for as
  boundary flux.
  \item The total charge in each sector is independently conserved.
\end{enumerate}
\end{theorem}

\begin{proof}[Proof of (a)]
For any interior vertex $v \in V_{L}$, all incident edges lie within $E_{L}$. The
LLC at $v$ in $G$ holds by hypothesis, and since all relevant edges are in $E_{L}$,
this is identical to the LLC for $G_{L}$. The argument is symmetric for $G_{R}$.
\end{proof}

\begin{proof}[Proof of (b)]
For a boundary vertex $v \in V_{L}$ with outgoing edges into $C$, the LLC at $v$
in $G$ states
\begin{equation}
  \sum_{\substack{u\in V_{L}\\(u,v)\in E_{L}}} \mathbf{q}(u)
  \;=\;
  \sum_{\substack{w\in V_{L}\\(v,w)\in E_{L}}} \mathbf{q}(w)
  \;+\;
  \sum_{(v,w)\in C} \mathbf{q}_{\mathrm{cut}}(w).
\end{equation}
This is the LLC at $v$ partitioned by destination---interior (in $E_{L}$) plus
boundary flux (in $C$).
\end{proof}

\begin{proof}[Proof of (c)]
Define the charge flow $f: E \to \mathbb{R}^{n}$ with the LLC as the discrete
divergence-free condition. By the discrete analogue of Gauss's theorem on $G$:
\begin{equation}
  \sum_{v \in V_{L}} \mathrm{div}(f)(v)
  \;=\;
  \sum_{e \in C} f(e)
  \;=\; \Phi(C).
\end{equation}
Since $\mathrm{div}(f)(v) = 0$ at all interior vertices (LLC holds), the left-hand
sum reduces to contributions from boundary vertices only. The total internal charge
is conserved: what enters $V_{L}$ at sources minus what leaves at sinks equals zero
by internal LLC, and the difference is exactly $\Phi(C)$. Therefore $Q_{L}$ is
conserved independently of $Q_{R}$.
\end{proof}

\subsubsection{Application: Event Horizons}

An event horizon in RAGC is a maximal causal cut $C^{*}$---one where every
future-directed causal path from the interior eventually intersects $C^{*}$. By
acyclicity of the DAG, this is well-defined.

\begin{corollary}[Horizon Partition]
Let $C^{*}$ be an event horizon cut. Then $G_{\mathrm{int}}$ and $G_{\mathrm{ext}}$
each satisfy the LLC independently. The charge in each sector is separately
conserved. The quantity $Q_{\mathrm{int}} + Q_{\mathrm{ext}}$ is not observable
from the exterior subgraph, because no edges cross from interior to exterior by
definition of a horizon. The question of a globally conserved total energy is
therefore not well-posed.
\end{corollary}

This result establishes the energy bookkeeping dissolution described in
Section~\ref{sec:severances}. As noted there, it does not resolve the unitarity
question; that requires the non-perturbative treatment of
Section~\ref{sec:hardwall}. A further scope note: the LLC result is a statement
about charge conservation in the discrete DAG graph. Its translation into the
language of full GR/QFT---where ``local conservation'' is $\nabla^{\mu}T_{\mu\nu}
= 0$ and depends on geometry, boundary terms, and the treatment of traced-out
degrees of freedom---requires the Bianchi compatibility work identified in
Section~\ref{sec:hardwall}(iv). The DAG Corollary is a necessary structural
result; showing it implies the correct GR boundary condition is an open task.

% ─── Summary ─────────────────────────────────────────────────────────────────

\subsection{Derivation 5: The covariant $\Box(\ln\lambda)$ field equation}
\label{sec:d5}

The RA Poisson equation established in \S\ref{sec:newtonian} gives
\begin{equation}
  \nabla^2(\ln\lambda) = -\frac{\rho_A}{3\xi}
  \label{eq:poisson_lnlam}
\end{equation}
in the quasi-static weak-field limit.  We now derive its fully covariant
generalisation, identify the Bianchi-compatible stress-energy tensor for the
$\lambda$-field, and discuss the new physics in the time-dependent and
strong-field regimes.

\subsubsection{Scalar-tensor action}

Define $\varphi \equiv \ln\lambda$.  The RA effective action for the
$\lambda$-field is
\begin{equation}
  S_\lambda[g,\varphi]
  = -\frac{\xi}{2}\int g^{\mu\nu}\nabla_\mu\varphi\,\nabla_\nu\varphi\,
      \sqrt{-g}\,\mathrm{d}^4x
  \;+\; \xi\int \bigl(g^{\mu\nu}\Box - \nabla^\mu\nabla^\nu\bigr)\varphi\;
      \delta g_{\mu\nu}\,\sqrt{-g}\,\mathrm{d}^4x \,.
  \label{eq:action_lambda}
\end{equation}
The second term is a non-minimal coupling (NMC) required for
Bianchi compatibility (\S\ref{sec:bianchi} below).

\subsubsection{Covariant field equation}

Taking the Euler--Lagrange variation of $S_\lambda + S_\mathrm{matter}$
with respect to $\varphi$:
\begin{equation}
  \xi\,\Box(\ln\lambda) = \frac{\partial\mathcal{L}_\mathrm{matter}}{\partial(\ln\lambda)}
    = -\frac{\rho_A}{3} \,,
\end{equation}
where the matter--$\lambda$ coupling $\partial\mathcal{L}/\partial(\ln\lambda)
= -\rho_A/3$ follows from the dimensional reduction structure of
$\mathcal{P}_\mathrm{act}$ in the three spatial directions (RAGC~\S4).
This gives
\begin{equation}
  \boxed{
    \Box(\ln\lambda)
    \;\equiv\; g^{\mu\nu}\nabla_\mu\nabla_\nu(\ln\lambda)
    \;=\; -\frac{\rho_A}{3\xi}
  }
  \label{eq:D5}
\end{equation}
as the covariant $\lambda$-field equation.

\subsubsection{Bianchi-compatible stress-energy tensor}
\label{sec:bianchi}

Varying the full action with respect to $g^{\mu\nu}$ yields the
$\lambda$-field stress-energy tensor:
\begin{equation}
  T_{\mu\nu}^\lambda
  = \xi\!\left[\nabla_\mu\varphi\,\nabla_\nu\varphi
      - \tfrac{1}{2}g_{\mu\nu}(\nabla\varphi)^2\right]
  + \xi\!\left[g_{\mu\nu}\Box\varphi - \nabla_\mu\nabla_\nu\varphi\right].
  \label{eq:Tmunu_lambda}
\end{equation}
The second line (the NMC term) is essential.  Taking the covariant
divergence and applying $[\Box,\nabla_\nu]\varphi = R_{\nu\mu}\nabla^\mu\varphi$:
\begin{equation}
  \nabla^\mu T_{\mu\nu}^\lambda
  = \xi\,\Box\varphi\cdot\nabla_\nu\varphi
    + \xi\,R_{\nu\mu}\nabla^\mu\varphi \,.
\end{equation}
Substituting~\eqref{eq:D5}: the first term becomes
$-(\rho_A/3)\nabla_\nu\varphi$, and the Ricci back-reaction term
$\xi R_{\nu\mu}\nabla^\mu\varphi$ encodes the covariant coupling
between spacetime curvature and $\lambda$-field gradients.
The full RA field equations
\begin{equation}
  G_{\mu\nu} = 8\pi G\,\mathcal{P}_\mathrm{act}[T_{\mu\nu}]
             + 8\pi G\,T_{\mu\nu}^\lambda
  \label{eq:RA_field_eqs}
\end{equation}
satisfy $\nabla^\mu G_{\mu\nu} = 0$ identically by the Bianchi identity,
provided~\eqref{eq:D5} holds.

\subsubsection{Newtonian and Hubble limits}

\paragraph{Quasi-static weak field.}
Setting $\partial_t = 0$ in~\eqref{eq:D5} with $(-+++)$ signature
gives $\Box \approx \nabla^2$, recovering~\eqref{eq:poisson_lnlam}. $\checkmark$

\paragraph{Vacuum propagation.}
In the absence of matter ($\rho_A = 0$):
\begin{equation}
  \Box(\ln\lambda) = 0 \,,
\end{equation}
so actualization-density fluctuations propagate as a massless scalar wave
at speed $c$.  This gives the $\lambda$-field a causal, retarded Green's
function — variations in one region influence $\lambda$ elsewhere only
after light-travel time.

\paragraph{Hubble tension.}
In the KBC void ($\rho_A \approx 0.7\,\rho_{A,\mathrm{global}}$),
$\Box(\ln\lambda) \approx 0$ in the void interior and
$\ln\lambda_{\mathrm{void}} > \ln\lambda_{\mathrm{global}}$.
The local Hubble rate satisfies $H_{\mathrm{local}} \propto \lambda_{\mathrm{local}}$
(RAGC~\S5.2), giving $H_{\mathrm{local}} > H_{\mathrm{CMB}}$ in agreement
with the observed tension.  The Eridanus supervoid prediction
($H \approx 75.9\;\mathrm{km\,s}^{-1}\mathrm{Mpc}^{-1}$,
$\delta \approx -0.30$) follows from the solution to~\eqref{eq:D5}
with the observed void density profile.

\subsubsection{Remaining open derivations}

Equation~\eqref{eq:D5} is the covariant field equation for $\ln\lambda$.
What remains is to prove:
\begin{enumerate}
  \item \emph{Bianchi compatibility in full non-linear GR}: the derivation above
    is perturbative (linear in $\lambda - \lambda_0$).  The non-linear
    proof requires showing $\nabla^\mu G_{\mu\nu} = 0$ holds for the full
    coupled system~\eqref{eq:RA_field_eqs}--\eqref{eq:D5}.
  \item \emph{Cosmological perturbation theory}: the coupled system of
    $\delta\lambda$, $\delta g_{\mu\nu}$ perturbations in an FRW background,
    needed for CMB power spectrum predictions.
  \item \emph{Strong-field corrections}: near black holes the field equation
    acquires additional terms from $R_{\nu\mu}\nabla^\mu\varphi$
    in~\eqref{eq:Tmunu_lambda}.
\end{enumerate}
Items 1--2 are targets for a modified-gravity collaborator
(Loll/Radboud, Deffner/UMBC).

\subsection{Derivation 6: Equivalence Principle Recovery and Operationalization
of $\lambda$}
\label{sec:d6}

\paragraph{Goal.}
Two foundational issues require explicit treatment before the dark matter
predictions (Derivation~5) can be considered publication-ready: (i) the recovery
of the weak equivalence principle (WEP) for ordinary matter from the
actualization-dependent sourcing equation, and (ii) the operational definition of
$\lambda(x)$ in terms of observables that would allow the Hubble tension prediction
(Prediction~1) to be tested quantitatively.

\paragraph{The equivalence principle tension.}
The WIMP prohibition rests on the claim that species with negligible
electromagnetic and strong cross-sections contribute negligible actualization flux
and therefore negligible metric sourcing. A referee correctly identifies that this
appears to imply gravitational coupling is not universal: it depends on
non-gravitational interaction channels, which conflicts with precision tests of the
weak equivalence principle (WEP) unless the framework provides an explicit
mechanism for universality recovery.

The resolution rests on the $\rho_A$ term of the unified sourcing
the unified sourcing equation $\nabla^2\Phi = 4\pi G[\rho_A + \rho_\lambda]$ (developed in RADM~\citep{Sandeman2026RADM}). For any species in thermodynamic equilibrium
over cosmological time, the accumulated causal depth satisfies
\begin{equation}
  \frac{A_{\mathrm{RA}}(x)}{V} \;\approx\; \lambda(x) \cdot t_{\mathrm{cosm}},
  \label{eq:equilibrium_limit}
\end{equation}
where $t_{\mathrm{cosm}}$ is the cosmological age of the system. In this limit,
$\rho_A \propto \lambda \cdot t_{\mathrm{cosm}} \cdot \langle\Delta E\rangle/c^2
\propto \rho_{\mathrm{matter}}$: the accumulated depth is proportional to ordinary
mass-energy, and $\rho_\lambda = 0$ (no gradient). Standard GR is recovered, and
the WEP is respected for all matter in equilibrium---neutrons, protons, lead atoms,
dark baryons if any. Precision tests of the WEP (E\"otv\"os experiments, lunar
laser ranging, atom interferometry) probe ordinary matter in quasi-equilibrium
conditions; these are the regime where $\rho_A \propto \rho_{\mathrm{matter}}$
exactly. Predicted outcome of torsion-balance tests: standard GR to current
precision.

WIMPs violate the universality only in the following precise sense: they have
$\lambda_{\mathrm{WIMP}} \approx 0$ at all times (not merely in the current epoch),
so $A_{\mathrm{RA,WIMP}} \approx 0$ even over cosmological time. They have
never been in actualization equilibrium with the baryonic sector. They are not a
species that violates universality; they are a species that has never participated
in the actualization dynamics that generate the metric in the first place.

\paragraph{Hard wall: formal proof of WEP recovery.}
The argument above establishes the equilibrium limit qualitatively but does not
constitute a proof. A publishable version of the WIMP prohibition requires:
\begin{enumerate}[noitemsep]
  \item A formal derivation showing that $\rho_A \to \rho_{\mathrm{matter}}$ in
  the equilibrium limit, with explicit bounds on the correction term
  $|\rho_A - \rho_{\mathrm{matter}}|/\rho_{\mathrm{matter}}$ as a function of
  departure from equilibrium.
  \item An explicit prediction for fifth-force experiments and atom interferometry:
  any species with $\lambda \ll \lambda_{\mathrm{baryon}}$ but non-negligible rest
  mass (e.g., a hypothetical feebly-interacting massive particle that has been in
  equilibrium at early times) would show a suppressed gravitational coupling
  proportional to its accumulated causal depth relative to baryons. This is a
  falsifiable prediction distinguishing RA from standard gravity.
  \item A demonstration that the $\rho_\lambda$ gradient correction is negligible
  at the scales probed by precision WEP tests (laboratory scales, $r \ll r_d$),
  recovering standard GR universality in the weak-field, quasi-static regime.
\end{enumerate}
\textbf{Collaborator target:} experimental gravity, precision tests of the
equivalence principle.

\paragraph{Operationalization of $\lambda(x)$.}
Prediction~1 (Hubble tension) requires an observable proxy for $\lambda(x)$. The
natural candidates are:
\begin{itemize}[noitemsep]
  \item \textbf{Galaxy number density} $n_{\mathrm{gal}}(x)$ from large-scale
  structure surveys (DESI, SDSS, 2MPP). This is proportional to local stellar
  mass density, which is proportional to integrated star-formation history, which
  is correlated with long-term actualization flux. Proxy quality: high, but
  includes selection effects.
  \item \textbf{Star-formation rate density} $\dot{\rho}_*$ from infrared surveys.
  More directly proportional to current $\lambda$ in dense regions, but misses
  the accumulated depth of old stellar populations.
  \item \textbf{Baryon density} $\rho_b$ from X-ray surveys and Sunyaev-Zel'dovich
  measurements. Proportional to total mass in a region, which correlates with
  long-term actualization history for matter in equilibrium.
\end{itemize}
The Hubble tension prediction becomes falsifiable precisely when: (a) $\lambda(x)$
is replaced by one of these observational proxies; (b) a smoothing scale is
specified (comoving $\sim 10$--$100$~Mpc to match BAO-scale structure); (c) the
predicted slope of $H_0$ residuals vs integrated density along the line of sight
is computed from the RA field equation; and (d) the result is compared against
existing $H_0$ maps from void-dominated surveys~\citep{Kenworthy2019} and
cluster-dominated surveys.

\paragraph{Hard wall: quantitative $H_0$ prediction.}
The current prediction is qualitatively distinguishable from $\Lambda$CDM and
early dark energy (which predict no line-of-sight correlation) but is not yet
a numerical forecast. The minimum requirement for a testable claim is:
\begin{enumerate}[noitemsep]
  \item Fix $\lambda \propto \rho_b$ (the simplest proxy) and derive the
  predicted $\delta H_0 / H_0$ as a function of the fractional overdensity
  $\delta$ along the line of sight.
  \item Show consistency with kinetic Sunyaev-Zel'dovich (kSZ) bounds on large
  void models, which already constrain strong line-of-sight $H_0$ variations.
  \item Produce a predicted slope and scatter for the $H_0$-density correlation
  using the RA field equation parameters.
\end{enumerate}

\paragraph{Upstream dependency: the generative measure $c_k$.}
The quantitative $H_0$ forecast is downstream of Derivation~1
(Section~\ref{sec:algebra}, Hard Wall~1). The bandwidth-partition formula
and Eridanus prediction are stated as Prediction~1b
(Section~\ref{sec:hubble_partition}). The ratio
$\bar\lambda_m/\bar\lambda_{\mathrm{exp}} \approx 0.42$ is currently calibrated
from the KBC void; it will be derived from the BDG-filter MCMC once Hard Wall~1
is resolved.
\textbf{Collaborator target:} observational cosmology, large-scale structure;
upstream: BDG-filter MCMC (D1).

\subsection{Summary: Status of Each Derivation}

\begin{description}
  \item[Derivation 1 (Generative Algebra).] \textit{Skeleton established.} The
  RA-CSG extension with on-shell filter is defined; causal invariance is inherited;
  process algebra formulation given. Mean-field Poisson-CSG formula derived:
  $c_k^{\mathrm{RA}}(x) = e^{-\mu}\mu^k/k!$ with $\mu = \lambda V_{\mathrm{coh}}
  p_{\mathrm{th}}$ (local, environment-dependent). BDG sign criterion
  (equation~\eqref{eq:BDG_filter}) identified as the discrete RA actualization
  filter. Hard wall: determine $V_{\mathrm{coh}}$ and run BDG-filter MCMC to
  extract asymptotic $c_k$ and $\bar\lambda_m/\bar\lambda_{\mathrm{exp}}$.

  \item[Derivation 2 (Macroscopic Limit and GR Inevitability).]
  \textit{Substantially resolved; restructured.} This derivation decomposes
  into three steps with distinct epistemic status.

  \textbf{D2a --- Lorentz invariance of the flat-space limit.} In the dense
  regime ($\mu \gg 1$, $\langle S_{\mathrm{BDG}}\rangle = 0$), the Poisson-CSG
  reduces to a Poisson process on Minkowski space with Lorentz-invariant
  statistics. The emergent metric is Minkowski. Substantially established by
  the Benincasa-Dowker conjecture~\citep{Benincasa2010} with strong analytical
  and numerical support. \textbf{Status: substantially resolved.}

  \textbf{D2b --- SR kinematics from discrete counting.} Given D2a, SR
  kinematics follow from integer counting on the causal graph
  (Section~\ref{subsec:bandwidth}): $\tau = \mathcal{N}(\gamma)\cdot t_P$
  (exact); $E = \gamma mc^2$ from the on-shell condition and Lorentz
  transformation of the actualization frequency; $c = l_P/t_P$ as maximum
  causal propagation rate. These are exact discrete results.
  \textbf{Status: resolved.}

  \textbf{D2c --- GR dynamics by Lovelock inevitability.} The Einstein
  field equations are the \emph{unique} consistent macroscopic description
  of the RA-CSG. By Lovelock's theorem~\citep{Lovelock1971}: in 4D, the
  only symmetric divergence-free rank-2 tensor from the metric and its
  derivatives is $G_{\mu\nu} + \Lambda g_{\mu\nu}$. The RA-CSG satisfies all
  three conditions: (i)~conservation --- LLC, Lean-verified; (ii)~covariance
  --- causal invariance, Lean-verified; (iii)~locality --- definitional in CSG.
  GR is not a limit to be recovered; it is logically guaranteed. The
  interesting physics is where RA \emph{departs} from GR: sparse regions
  ($\mu \sim 1$), Planck-scale corrections, discrete singularity structure,
  gravitational-wave dispersion. \textbf{Status: resolved by Lovelock.
  Planck-scale corrections: open.}

  \textbf{Remaining genuine gap (D2c$'$):} Extension of frame-independence
  to the type~III$_1$ AQFT setting via Araki's modular relative
  entropy~\citep{Araki1976}: proving $\Delta S_{\mathrm{Araki}}(\rho_{\mathrm{KMS}}
  \| \Omega) = 0$ under modular flow and $\Pact$ commutes with
  $\sigma_s^{\Omega}$. Needed for the AQFT Unruh extension, not for GR.
  \textbf{Collaborator target:} algebraic QFT, Tomita-Takesaki modular theory.

  \item[Derivation 3 ($\Pact$ Type-Constraint).] \textit{Substantial partial result.}
  Spectral definition given; vacuum energy suppression proved
  (Theorem~\ref{thm:vacuum}); gauge covariance established (Proposition~2); DAG
  connection formalized. Hard walls (four): non-perturbative regime, renormalization
  compatibility (including treatment of Casimir/binding energy corrections),
  back-reaction fixed point, and---most fundamentally---Bianchi compatibility
  ($\nabla^{\mu}\langle\Tact\rangle = 0$). The recommended path forward is to
  reframe $\Pact$ as a covariant renormalization prescription rather than a
  literal projector that excludes virtual contributions from gravity.


% ─── Thermodynamic Derivation (inserted after D3) ────────────────────────────
\subsection{Thermodynamic Derivation of the Actualization Tensor}
\label{sec:thermo_derivation}

While the 2D cylindrical gradient correction
$\Sigma_\lambda = \xi/(4\pi G R r_d)$ gives flat galactic rotation curves
in the disk regime (RADM~\S3.1, RA-native), and the DJW sparse-halo regime
gives the lensing convergence $\kappa(b) = v_\mathrm{flat}^2/(4G\Sigma_\mathrm{cr}b)$
as a theorem of the causal structure (RADM~\S4.1), establishing the covariant
origin of the $\lambda$-field sourcing term is paramount. Building on Jacobson's thermodynamic derivation of the Einstein
equation of state~\citep{Jacobson1995}, we map the local Rindler horizon to the
RA causal severance boundary---the topological limit across which actualization
directed edges have not yet formed.

The fundamental Clausius relation across this boundary is $dQ = T\,dS$. In
standard general relativity the entropy variation $dS$ is sourced entirely by the
heat flux of accumulated matter, $dS_A$. However, because Relational Actualism
defines the kinematic actualization vertex strictly as an irreversible increase in
relative entanglement entropy ($\Delta S(\rho\|\sigma_0) > 0$), the background
generation of the causal DAG itself constitutes a fundamental entropy flux. We
therefore decompose the total entropy variation:
\begin{equation}
  dS_{\mathrm{total}} = dS_A + dS_\lambda.
\end{equation}

To formulate $dS_\lambda$ covariantly, we treat the logarithmic actualization
density $\phi \equiv \ln\lambda$ as a scalar field. A constant, uniform
actualization flow ($\nabla_\mu\phi = 0$) represents a steady-state vacuum,
yielding no net horizon deformation. A spatial gradient in actualization density
necessitates an acceleration of entropy production, requiring second derivatives
of $\phi$. The unique, dimensionally consistent tensor coupling this scalar to
the horizon geometry is:
\begin{equation}
  \Theta_{\mu\nu}^{(\lambda)} = \xi
    \Bigl[\nabla_\mu\nabla_\nu(\ln\lambda)
          - g_{\mu\nu}\,\Box(\ln\lambda)\Bigr].
  \label{eq:actualization_tensor}
\end{equation}
The actualization entropy flux across the horizon, generated by the null
congruence $k^\mu$, is thus
$T\,dS_\lambda = \int_{\mathcal{H}} \Theta_{\mu\nu}^{(\lambda)}k^\mu k^\nu\,d\Sigma$.
Because $k^\mu$ is null ($g_{\mu\nu}k^\mu k^\nu = 0$), the d'Alembertian term
drops out of the surface integral, strictly mirroring the behaviour of the
cosmological constant in Jacobson's original formulation.

Equating the expansion of the horizon (via the Raychaudhuri equation) with the
total entropy flux yields the RA-modified field equations:
\begin{equation}
  G_{\mu\nu} = 8\pi G\!\left(T_{\mu\nu}^{(A)}
    + \Theta_{\mu\nu}^{(\lambda)}\right).
  \label{eq:RA_field_eqs}
\end{equation}
In the static weak-field limit the $00$-component of this tensor recovers
exactly the $\nabla^2(\ln\lambda)$ Poisson modification proposed in RADM.
We now give the explicit derivation.

\begin{proposition}[Newtonian limit of the covariant actualization term]
\label{prop:newtonian_limit}
Let $\phi \equiv \ln\lambda$ be static ($\partial_t\phi = 0$) and let the
spacetime metric be the isotropic weak-field metric
\begin{equation}
  ds^2 = -(1+2\Phi)\,dt^2 + (1-2\Phi)\,\delta_{ij}\,dx^i\,dx^j,
  \quad |\Phi| \ll 1.
\end{equation}
Then
\begin{equation}
  \Box(\ln\lambda) \;=\; \nabla^2(\ln\lambda) + O\!\left(\Phi\,|\nabla^2(\ln\lambda)|\right).
  \label{eq:box_limit}
\end{equation}
Equivalently, the correction is suppressed by the gravitational potential
$\Phi \sim 10^{-6}$ (galactic) to $10^{-9}$ (solar system), negligible
for all current observational tests.
\end{proposition}

\begin{proof}
Write $\Box\phi = g^{\mu\nu}\nabla_\mu\nabla_\nu\phi
= g^{\mu\nu}(\partial_\mu\partial_\nu\phi - \Gamma^\rho_{\mu\nu}\partial_\rho\phi)$.
Since $\phi$ is static the time component $g^{00}\partial_0\partial_0\phi = 0$.
The spatial kinetic term gives
\[
  g^{ij}\partial_i\partial_j\phi \;=\; (1+2\Phi)\,\delta^{ij}\partial_i\partial_j\phi
  \;=\; \nabla^2\phi + O(\Phi\,|\nabla^2\phi|).
\]
For the Christoffel correction $-g^{\mu\nu}\Gamma^\rho_{\mu\nu}\partial_\rho\phi$,
the two nonzero contributions to first order in $\Phi$ are:
\begin{enumerate}[noitemsep]
  \item \emph{Temporal sector} ($\mu=\nu=0$, $\rho=i$):
  $\Gamma^i_{00} = \partial^i\Phi + O(\Phi^2)$, so
  $-g^{00}\Gamma^i_{00}\partial_i\phi = +(1-2\Phi)\,(\nabla\Phi)\cdot(\nabla\phi)
  = +(\nabla\Phi)\cdot(\nabla\phi) + O(\Phi^2)$.
  \item \emph{Spatial sector} ($\mu=i$, $\nu=j$, $\rho=k$):
  The trace of the spatial Christoffel symbol is
  $\delta^{ij}\Gamma^k_{ij} = +\partial^k\Phi + O(\Phi^2)$
  (using $\Gamma^k_{ij} = -({\delta^k_j\partial_i\Phi + \delta^k_i\partial_j\Phi
  - \delta_{ij}\partial^k\Phi})$), so
  $-g^{ij}\Gamma^k_{ij}\partial_k\phi = -(1+2\Phi)\,(\nabla\Phi)\cdot(\nabla\phi)
  = -(\nabla\Phi)\cdot(\nabla\phi) + O(\Phi^2)$.
\end{enumerate}
The two $(\nabla\Phi)\cdot(\nabla\phi)$ contributions cancel exactly at $O(\Phi)$,
leaving $\Box\phi = \nabla^2\phi + O(\Phi|\nabla^2\phi|)$.

\smallskip\noindent
\emph{Remark.} The cancellation is structural: the isotropic-gauge metric is
conformally flat to $O(\Phi)$, so the temporal and spatial Christoffel traces
are equal in magnitude and opposite in sign. It holds for any scalar field
in any conformally flat perturbation of Minkowski space.
\end{proof}

\begin{corollary}[Newtonian limit of the RA field equations]
\label{cor:poisson_ra}
In the static weak-field limit, the $00$-component of
equation~\eqref{eq:RA_field_eqs} yields the modified Poisson equation
\begin{equation}
  \nabla^2\Phi \;=\; 4\pi G\,\rho_A \;+\; \xi\,\nabla^2(\ln\lambda),
  \label{eq:modified_poisson}
\end{equation}
which recovers the unified sourcing equation of
RADM~\citep{Sandeman2026RADM} with the identification
$\rho_\lambda \equiv \xi/(4\pi G)\cdot\nabla^2(\ln\lambda)$.
\end{corollary}

\begin{proof}
The $00$-component of $\Theta_{\mu\nu}^{(\lambda)}$ is
$\Theta_{00} = \xi[\nabla_0\nabla_0\phi - g_{00}\,\Box\phi]$.
Since $\phi$ is static, $\nabla_0\nabla_0\phi = \partial_t^2\phi = 0$.
Using $g_{00} = -(1+2\Phi)$ and Proposition~\ref{prop:newtonian_limit}:
$\Theta_{00} = \xi(1+2\Phi)\nabla^2(\ln\lambda) = \xi\,\nabla^2(\ln\lambda) + O(\Phi)$.
Inserting into the field equations and matching the linearised Einstein tensor
$G_{00} = 2\nabla^2\Phi$ yields~\eqref{eq:modified_poisson}.
\end{proof}

\noindent
The coupling $\xi$ is fixed by the microscopic RA entropy
(equation~\eqref{eq:xi_microscopic}); it is not a free parameter.

Furthermore, by strict analogy with the Bekenstein-Hawking area-entropy
relation, the coupling constant $\xi$ is not a free parameter. It must emerge
as the ratio of the discrete entropy generated per actualization vertex to the
fundamental RA area element:
\begin{equation}
  \xi \;\propto\; \frac{\langle\Delta S_{\mathrm{vertex}}\rangle}{l_{\mathrm{RA}}^2}.
  \label{eq:xi_microscopic}
\end{equation}
Crucially, this connects macroscopic gravity directly to the Kinematic Coherence
Bound established in RAQI~\citep{Sandeman2026RAQI}. The rate limit on
actualization events per coherence volume pins this entropy ratio to an
observable threshold, constraining $\xi$ via fault-tolerant quantum array
decoherence limits.

\subsubsection*{The Actualization Field Equation and Its Observational Consequences}

\paragraph{The open question is not Bianchi — it is $\xi$.}
The contracted Bianchi identity $\nabla^\mu G_{\mu\nu} = 0$ is a property of
the emergent continuum metric; RA derives this metric from the causal graph via
the Jacobson thermodynamic construction. The LLC is the RA conservation law
(Lean-verified), and Bianchi is what the LLC becomes in the continuum limit.
Framing this as a gap would have the logical arrow backwards.

The genuinely open question is not whether RA is self-consistent in the continuum
limit --- it is --- but what the coupling constant $\xi$ is, and what the
actualization density $\lambda(x)$ predicts quantitatively.

\paragraph{The $\lambda$ equation of motion.}
Requiring consistency of the RA field equations~\eqref{eq:RA_field_eqs} with
the covariant divergence condition yields a derived equation of motion for the
actualization field:
\begin{equation}
  \Box(\ln\lambda) \;=\; -\frac{\rho_A}{3\xi},
  \label{eq:lambda_eom}
\end{equation}
where $T = -\rho_A$ for pressureless dust. In the static weak-field limit
(Proposition~\ref{prop:newtonian_limit}):
\begin{equation}
  \nabla^2(\ln\lambda) \;=\; -\frac{\rho_A}{3\xi}.
  \label{eq:lambda_poisson}
\end{equation}
This is a \emph{prediction} of RA, not a constraint imposed from outside.
It says: $\lambda$ is sourced by baryonic matter density, just as $\Phi$
is sourced by $\rho_A$ in the Poisson equation. The actualization density
field is not freely specifiable --- it is determined by the matter distribution.

\paragraph{Physical content of~\eqref{eq:lambda_poisson}.}
Equation~\eqref{eq:lambda_poisson} predicts a specific spatial profile for $\lambda$:
maximised at baryonic concentrations, falling off toward galactic outskirts.
The dark-matter contribution $\rho_\lambda \propto \nabla^2(\ln\lambda)$ is therefore
concentrated at the baryonic-to-void \emph{transition} --- the galactic halo region.
This is structurally consistent with flat rotation curves and Bullet Cluster dynamics
(RADM~\citep{Sandeman2026RADM}), and provides a new RA-native derivation of
that spatial profile directly from the field equations rather than by postulate.

\paragraph{Screening and solar system consistency.}
The effective mass of $\ln\lambda$ from~\eqref{eq:lambda_poisson} is
$m_{\mathrm{eff}}^2 \approx \rho_A/(3\xi)$, which scales with local matter density.
In the solar neighbourhood $m_{\mathrm{eff}}$ is large, giving Yukawa suppression
$e^{-m_{\mathrm{eff}} r}$ on AU scales. The fifth-force range at galactic scales
is $r_d \sim (3\xi/\bar\rho_A)^{1/2}$, which is a new RA prediction: the
dark-matter halo scale-length is set by the ratio of $\xi$ to the mean baryonic
density.

\paragraph{Open target: fixing $\xi$ from RA.}
The coupling constant $\xi$ appears in all dark-matter predictions but is not
yet derived from RA's microscopic structure. From~\eqref{eq:xi_microscopic},
$\xi \propto \langle\Delta S_{\mathrm{vertex}}\rangle / l_{\mathrm{RA}}^2$, which
connects $\xi$ to the entropy generated per actualization event and the
Planck-scale RA area element. The Kinematic Coherence Bound of
RAQI~\citep{Sandeman2026RAQI} pins the actualization rate per coherence volume,
providing a second handle. The derivation of $\xi$ from these two inputs is
the primary RA-internal open target for this derivation.
\textbf{Collaborator target:} quantum information, quantum gravity
phenomenology, large-scale structure.
  \item[Derivation 4 (Graph Cut Theorem).] \textit{Complete at the DAG level;
  \textbf{fully verified in Lean~4/Mathlib with no \texttt{sorry} tags}.}
  Full proof of Theorem~\ref{thm:cut} (Ledger Partition) with three parts; horizon
  corollary derived. The discrete graph-theoretic core has been verified in the
  Lean~4 theorem prover, including full proofs of the two sum-decomposition helper
  lemmas (\texttt{sum\_outgoing\_decompose} and \texttt{sum\_incoming\_decompose}).
  The Lean file is completely \texttt{sorry}-free for all Graph Cut results.
  Scope: energy bookkeeping in the discrete DAG, not unitarity. Translation to
  full GR boundary conditions requires Bianchi compatibility (Derivation~3,
  hard wall~iv).

  \item[Derivation 7 (Covariant Step 4 --- New, from Engine of Becoming).]
  \textit{Open.} The stochastic graph update in the Engine of
  Becoming~\citep{Sandeman2026RAEB} selects one actualization vertex per
  algorithmic step --- a description that is foliation-dependent. A covariant
  formulation must show that the stochastic update commutes for
  spacelike-separated candidates: selecting $v_A$ then $v_B$ produces the same
  graph as selecting $v_B$ then $v_A$ when $v_A$ and $v_B$ are causally
  disconnected. This is the Tomonaga-Schwinger analogue for the stochastic
  update, and it is the condition required for the full covariant quantum
  gravity interpretation of the algorithm. The causal invariance already
  established in Derivation~1 (for the deterministic graph growth) provides the
  nearest existing scaffolding.
  \textbf{Hard wall:} proving the stochastic update is causally invariant for
  spacelike-separated candidates, with an explicit error bound.
  \textbf{Collaborator target:} algebraic QFT, causal set theory.

  \item[Derivation 5 (Unified Metric Sourcing).] \textit{Relocated to RADM.}
  The $\nabla^2(\ln\lambda)$ conjectural mechanism, rotation curve derivation,
  Bullet Cluster $A_{\mathrm{RA}}$ argument, Tully-Fisher scaling, and CMB
  structure formation requirements are the primary subject of the companion paper
  RADM~\citep{Sandeman2026RADM}. The WIMP prohibition in Prediction~4 stands
  independently and is supported by the equilibrium-limit WEP recovery argument
  in Derivation~6.

  \textit{Substantially complete.} The covariant d'Alembertian
  $\Box(\ln\lambda) = g^{\mu\nu}\nabla_\mu\nabla_\nu(\ln\lambda)$ is
  the correct covariant extension.
  \textbf{Proved} (Proposition~\ref{prop:newtonian_limit},
  Corollary~\ref{cor:poisson_ra}): $\Box(\ln\lambda) = \nabla^2(\ln\lambda)
  + O(\Phi)$ in the static weak-field limit; the modified Poisson equation
  $\nabla^2\Phi = 4\pi G\rho_A + \xi\,\nabla^2(\ln\lambda)$ is derived.
  \textbf{Derived} (equation~\eqref{eq:lambda_poisson}): the actualization
  density satisfies $\nabla^2(\ln\lambda) = -\rho_A/(3\xi)$, sourced by
  baryonic matter. This predicts the dark-matter halo profile from
  first principles.
  \textbf{Primary open target:} fixing $\xi$ from RA microscopic
  structure (entropy per vertex and RAQI Kinematic Coherence Bound).
  \textbf{Collaborator targets:} quantum gravity phenomenology,
  large-scale structure, precision gravity.

  \item[Derivation 6 ($\xi$ Determination and Quantitative $\lambda$ Predictions).]
  \textit{Framework established; quantitative predictions outstanding.}
  WEP recovery for ordinary matter established via equilibrium limit
  ($\rho_A \to \rho_{\mathrm{matter}}$); WIMP prohibition maintained
  ($A_{\mathrm{RA}} \approx 0$ for zero-cross-section species).
  The $\lambda$ equation of motion $\nabla^2(\ln\lambda) = -\rho_A/(3\xi)$
  is derived. Observable proxies for $\lambda(x)$ identified:
  galaxy density, star-formation rate, baryon density.
  \textbf{Open targets:} (i) derive $\xi$ from
  $\langle\Delta S_{\mathrm{vertex}}\rangle/l_{\mathrm{RA}}^2$ and the
  RAQI Kinematic Coherence Bound; (ii) predict the $H_0$-density slope
  quantitatively from $\xi$ and $\bar\rho_A$; (iii) predict the
  Tully-Fisher normalisation from $\xi$; (iv) kSZ velocity field from
  $\lambda$ profile; (v) fifth-force range $r_d = \sqrt{3\xi/\bar\rho_A}$
  as a falsifiable prediction of RA.
  \textbf{Collaborator targets:} quantum gravity phenomenology,
  observational cosmology, large-scale structure surveys.
\end{description}

% ─────────────────────────────────────────────────────────────────────────────

% ─────────────────────────────────────────────────────────────────────────────
\section*{Epilogue: The Horizon of the Program --- From Gravity to Gauge Symmetries}
\addcontentsline{toc}{section}{Epilogue: The Horizon of the Program}
\label{sec:epilogue}
% ─────────────────────────────────────────────────────────────────────────────

Relational Actualism strips gravity of its status as a fundamental force and
reconceives it as the macroscopic thermodynamic shadow of the causal graph's
accumulated topology. This suggests---if the framework is correct---that
quantizing the metric is the wrong approach to quantum gravity. What requires
understanding is not a quantum field on a spacetime background, but the discrete
actualization events that generate the causal graph from which spacetime emerges.

In its current formulation, however, the framework is agnostic about the origins
of the remaining three fundamental forces and the specific particle content of the
Standard Model. RA currently borrows the Standard Model as a phenomenological
engine: it uses known electromagnetic, weak, and strong cross-sections to inform
the graph's transition probabilities without deriving those cross-sections from
the graph structure itself. A mature theory of fundamental physics cannot borrow
its interaction rules; it must derive them.

\paragraph{Progress on this horizon.}
The companion paper RASM~\citep{Sandeman2026RASM} has now substantially
advanced this programme.  The BDG integers $(1,-1,9,-16,8)$ determine a
complete, Lean-verified five-element particle topology (RATM), give exactly
three non-gravitational gauge interactions, and yield two proved coupling
constants:
\begin{equation}
  \alpha_s(m_Z) = \frac{1}{\sqrt{c_2 c_4}} = \frac{1}{\sqrt{72}} = 0.11785
  \qquad (\text{PDG: } 0.11800,\; 0.13\%)
\end{equation}
from the BDG path weight at $\mu=1$~\citep{Sandeman2026BDG}, alongside the
Wyler formula for $\alpha_{\mathrm{EM}}$ (verified to $0.0001\%$).  The
dark matter/baryon ratio follows from the same integers:
$\Omega_{\mathrm{CDM}}/\Omega_b = 5.40$ (Planck~2018: $5.416$, $0.3\%$,
Prediction~5 above).

The BDG Regge formula $m_n^2 = 4\sigma n$ (with $a_0=0$ from the identity
$c_1+c_2+c_3+c_4=0$) gives the proton mass:
\begin{equation}
  m_p = \sqrt{c_4}\,\Lambda_{\mathrm{QCD}} = \sqrt{8}\times 332~\mathrm{MeV}
  = 939~\mathrm{MeV} \qquad (0.08\%),
\end{equation}
and the Koide scale $m_0 = m_p/3 = \sqrt{c_4/c_2}\,\Lambda_{\mathrm{QCD}} = 313$~MeV ($0.08\%$),
connecting the cosmological derivation target $f_0$ to the hadronic string
tension $\sigma = \Lambda_{\mathrm{QCD}}^2$. The RAGC Derivation~1 simulation
therefore simultaneously predicts the cosmological dark matter density,
the proton mass, and the Koide scale with zero free parameters.

What was stated as a long-range research direction in the original RAGC
is now a partially completed derivation.

The ultimate multi-decade horizon of the RA program is the derivation of the
Standard Model from the pure discrete topology of the causal DAG. What follows
is a statement of the research direction, not a claim of achievement.

\paragraph{Particles as topological motifs.}
In a pure process ontology, what we call ``matter'' is the accumulated causal
structure of on-shell actualization events. If that is so, the specific
\emph{kinds} of matter cannot be ontologically primitive; they must be structural
features of the graph. The natural candidate is the \emph{topological motif}: a
localised, stable, repeating causal structure---a knot in the DAG---whose
characteristic actualization frequency determines its contribution to $A_{\mathrm{RA}}$
and thereby its effective inertial mass. Different particle species (an electron
versus an up quark) would correspond to mathematically distinct knotting patterns.
This direction has precedents in the physics literature: Bilson-Thompson's braid
model of preons~\citep{BilsonThompson2006}, Hackett and Smolin's work on particle
states in spin foam models, and Penrose's original spinor networks all explore
the idea that particle quantum numbers can be read off topological invariants of
discrete graph structures.

\paragraph{Gauge symmetries as edge-routing conservation laws.}
If particles are topological motifs, the forces between them must arise from the
rules governing how motifs interact---how causal knots exchange edges. The
continuous gauge symmetries of the Standard Model ($U(1)$, $SU(2)$, $SU(3)$)
are conjectured to emerge as macroscopic approximations of discrete edge-routing
and topological conservation laws on the DAG:

\begin{description}[leftmargin=1.5em]
  \item[\textbf{Electromagnetism ($U(1)$)}] The simplest Abelian case.
  Conjectured to emerge from a single discrete conservation law on directed
  edges---the graph-theoretic analogue of charge conservation---whose
  continuous limit is the $U(1)$ gauge symmetry of quantum electrodynamics.

  \item[\textbf{The Weak Force ($SU(2)$) and Chirality}] The most
  structurally puzzling of the three, because the weak interaction violates
  parity---it couples exclusively to left-handed fermions. In graph terms, this
  would require the DAG to have an intrinsic handedness: a direction convention
  on causal edges that distinguishes left-handed from right-handed motifs. How a
  directed acyclic graph acquires an intrinsic chirality is an open problem with
  no clear current path. This is the hardest of the three gauge symmetries to
  derive and must be named as such.

  \item[\textbf{The Strong Force ($SU(3)$) and Confinement}] Quarks are
  empirically never isolated (colour confinement). In the DAG picture, this
  suggests that certain topological motifs (quarks) are structurally unstable as
  open subgraphs: they do not achieve causal closure on their own. They must
  bind in specific multiplicities---triplets (baryons) or pairs (mesons)---to
  form a closed, stable causal subgraph. Colour confinement would then be a
  topological stability criterion rather than a dynamical force, an analogue of
  why the Markov blanket must be closed to maintain autopoiesis. The $SU(3)$
  group structure would emerge from the symmetries of the allowed closure
  patterns.
\end{description}

\paragraph{The hard wall.}
These conjectures are statements of a research direction, not derivations. The
hard wall is foundational: connecting the abstract topology of the causal DAG to
the specific group-theoretic structure of the Standard Model gauge group
$U(1) \times SU(2) \times SU(3)$ requires a mathematical bridge between discrete
graph theory and continuous Lie algebra representation theory. The most promising
existing frameworks for this bridge are spin foam models and loop quantum gravity
(which embed $SU(2)$ in the graph structure directly) and category-theoretic
approaches to topological quantum field theory. RA's contribution is to provide
a clearer physical motivation for why the derivation should be attempted: if
actualization events are ontologically primitive and gravity is their thermodynamic
shadow, the particle content of the universe is not a brute fact but a theorem
waiting to be proved.


\section*{Declarations}
\textbf{Funding:} Not applicable.

\textbf{Ethics statement:} This article does not contain any studies with human participants or animals performed by the author.

\quad
\textbf{Competing interests:} The author declares no competing interests.

\textbf{Preprint availability:} This manuscript is available as a preprint at
\href{https://doi.org/10.5281/zenodo.19197999}{doi.org/10.5281/zenodo.19197999}.
Companion papers: EB~\href{https://doi.org/10.5281/zenodo.19198120}{doi.org/10.5281/zenodo.19198120}, RADM~\href{https://doi.org/10.5281/zenodo.19198513}{doi.org/10.5281/zenodo.19198513}.

\textbf{AI use disclosure:} Claude (Anthropic, model \texttt{claude-sonnet-4-6}) was used in the preparation of this manuscript, including drafting and revision of text, mathematical derivation assistance, \LaTeX{} compilation, and Lean~4 proof development. The author is solely responsible for all scientific content, mathematical claims, and conclusions. Gemini (Google DeepMind, Gemini~2.0~Pro) was used for supplementary literature checking and cross-verification of symbolic calculations.


\textbf{Code availability:} The Lean~4 formal proof files (\texttt{RA\_AQFT\_Proofs\_v10.lean} and \texttt{RA\_Proofs\_Lean4.lean}), Python computational scripts, and all companion paper source files are publicly available at \href{https://github.com/jsandeman/Relational-Actualism}{github.com/jsandeman/Relational-Actualism}.

\textbf{Data availability:} This paper presents theoretical work. No datasets
were generated or analysed.

\bibliographystyle{plainnat}
\begin{thebibliography}{99}


\bibitem[Sandeman(2026a)]{Sandeman2026a}
J.~F.~Sandeman.
\newblock The Relational Universe: Interaction and the Transition of the
  Possible to the Actual ({RAQM}).
\newblock \textit{Foundations of Physics} (Under review), 2026.
\newblock Preprint: \href{https://doi.org/10.5281/zenodo.19174380}{doi.org/10.5281/zenodo.19174380}.

\bibitem[Sandeman(2026b)]{Sandeman2026b}
J.~F.~Sandeman.
\newblock Relational Actualism: Gravity and Cosmology ({RAGC}).
\newblock Preprint, 2026.
\newblock \href{https://doi.org/10.5281/zenodo.19197999}{doi.org/10.5281/zenodo.19197999}.

\bibitem[Sandeman(2026c)]{Sandeman2026c}
J.~F.~Sandeman.
\newblock Relational Actualism: Complexity and Intelligence ({RACI}).
\newblock Preprint, 2026.
\newblock \href{https://doi.org/10.5281/zenodo.19198224}{doi.org/10.5281/zenodo.19198224}.

\bibitem[Sandeman(2026d)]{Sandeman2026RAHC}
J.~F.~Sandeman.
\newblock Algorithmic Causal Dynamics: A Unified Framework for Emergence and
  Complexity in Relational Actualism ({RAHC}).
\newblock Preprint, 2026.
\newblock \href{https://doi.org/10.5281/zenodo.19198171}{doi.org/10.5281/zenodo.19198171}.

\bibitem[Sandeman(2026e)]{Sandeman2026RAEB}
J.~F.~Sandeman.
\newblock The Engine of Becoming: A Covariant Stochastic Graph Rewriting
  Algorithm as the Unified Foundation of Quantum Mechanics and General
  Relativity ({EB}).
\newblock Preprint, 2026.
\newblock \href{https://doi.org/10.5281/zenodo.19198120}{doi.org/10.5281/zenodo.19198120}.

\bibitem[Sandeman(2026f)]{Sandeman2026RAQI}
J.~F.~Sandeman.
\newblock Relational Actualism: Implications for Quantum Information ({RAQI}).
\newblock Preprint, 2026.
\newblock \href{https://doi.org/10.5281/zenodo.19198285}{doi.org/10.5281/zenodo.19198285}.

\bibitem[Sandeman(2026f)]{Sandeman2026RADM}
J.~F.~Sandeman.
\newblock Relational Actualism: Dark Matter as Actualization-Density Topology
  ({RADM}).
\newblock Preprint, 2026.
\newblock \href{https://doi.org/10.5281/zenodo.19198513}{doi.org/10.5281/zenodo.19198513}.

\bibitem[Rovelli(1996)]{Rovelli1996}
C.~Rovelli.
\newblock Relational quantum mechanics.
\newblock \textit{International Journal of Theoretical Physics},
  35(8):1637--1678, 1996.

\bibitem[Rideout and Sorkin(2000)]{RideoutSorkin2000}
D.~P.~Rideout and R.~D.~Sorkin.
\newblock Classical sequential growth dynamics for causal sets.
\newblock \textit{Physical Review D}, 61(2):024002, 2000.

\bibitem[Bland-Hawthorn \& Gerhard(2016)]{BlandHawthorn2016}
J.~Bland-Hawthorn and O.~Gerhard.
\newblock The Galaxy in context: Structural, kinematic, and integrated
  properties.
\newblock \emph{Annual Review of Astronomy and Astrophysics}, 54:529--596, 2016.
\newblock \href{https://doi.org/10.1146/annurev-astro-081915-023441}{doi:10.1146/annurev-astro-081915-023441}.

\bibitem[Benincasa \& Dowker(2010)]{Benincasa2010}
D.~M.~T.~Benincasa and F.~Dowker.
\newblock The scalar curvature of a causal set.
\newblock \emph{Physical Review Letters}, 104(18):181301, 2010.
\newblock \href{https://doi.org/10.1103/PhysRevLett.104.181301}{doi:10.1103/PhysRevLett.104.181301}.

\bibitem[Durhuus et~al.(2010)]{Durhuus2010}
B.~Durhuus, T.~Jonsson, and J.~F.~Wheater.
\newblock On the spectral dimension of causal triangulations.
\newblock \emph{Journal of Statistical Physics}, 139:859--881, 2010.
\newblock \href{https://doi.org/10.1007/s10955-010-9968-x}{doi:10.1007/s10955-010-9968-x}.

\bibitem[Sorkin(2003)]{Sorkin2003}
R.~D.~Sorkin.
\newblock Causal sets: Discrete gravity.
\newblock In A.~Gomberoff and D.~Marolf, editors, \textit{Lectures on Quantum
  Gravity}, NATO Science Series, pages 305--327. Springer, 2003.

\bibitem[Dowker(2005)]{Dowker2005}
F.~Dowker.
\newblock Causal sets and the deep structure of spacetime.
\newblock In A.~Ashtekar, editor, \textit{100 Years of Relativity}, pages
  445--467. World Scientific, 2005.

\bibitem[Bombelli et~al.(2009)]{Bombelli2009}
L.~Bombelli, J.~Henson, and R.~D.~Sorkin.
\newblock Discreteness without symmetry breaking: a theorem.
\newblock \textit{Modern Physics Letters A}, 24(32):2579--2587, 2009.

\bibitem[Ambjørn et~al.(2012)]{Ambjorn2012}
J.~Ambjørn, J.~Jurkiewicz, and R.~Loll.
\newblock Causal dynamical triangulations and the quest for quantum gravity.
\newblock In G.~Ellis, J.~Murugan, and A.~Weltman, editors,
  \textit{Foundations of Space and Time}, pages 321--337.
  Cambridge University Press, 2012.

\bibitem[Weinberg(1989)]{Weinberg1989}
S.~Weinberg.
\newblock The cosmological constant problem.
\newblock \textit{Reviews of Modern Physics}, 61(1):1--23, 1989.

\bibitem[Kaloper and Padilla(2014)]{KaloperPadilla2014}
N.~Kaloper and A.~Padilla.
\newblock Sequestering the standard model vacuum energy.
\newblock \textit{Physical Review Letters}, 112(9):091304, 2014.

\bibitem[Bassi et~al.(2013)]{Bassi2013}
A.~Bassi, K.~Lochan, S.~Satin, T.~P.~Singh, and H.~Ulbricht.
\newblock Models of wave-function collapse, underlying theories, and
  experimental tests.
\newblock \textit{Reviews of Modern Physics}, 85(2):471--527, 2013.

\bibitem[Bisognano and Wichmann(1975)]{BisognanoWichmann1975}
J.~J.~Bisognano and E.~H.~Wichmann.
\newblock On the duality condition for a Hermitian scalar field.
\newblock \textit{Journal of Mathematical Physics}, 16(5):985--1007, 1975.

\bibitem[Araki(1976)]{Araki1976}
H.~Araki.
\newblock Relative entropy of states of von {Neumann} algebras.
\newblock \textit{Publications of the Research Institute for Mathematical
  Sciences}, 11(3):809--833, 1976.

\bibitem[Jacobson(1995)]{Jacobson1995}
T.~Jacobson.
\newblock Thermodynamics of spacetime: the {Einstein} equation of state.
\newblock \textit{Physical Review Letters}, 75(7):1260--1263, 1995.

\bibitem[Wald(1994)]{Wald1994}
R.~M.~Wald.
\newblock \textit{Quantum Field Theory in Curved Spacetime and Black Hole
  Thermodynamics}.
\newblock University of Chicago Press, Chicago, 1994.

\bibitem[Sulis(2020)]{Sulis2020}
W.~Sulis.
\newblock Process algebra models of quantum mechanics.
\newblock \textit{Physics Essays}, 33(3):333--340, 2020.

\bibitem[Bilson-Thompson(2006)]{BilsonThompson2006}
S.~O.~Bilson-Thompson.
\newblock A topological model of composite preons.
\newblock arXiv:hep-ph/0503213, 2006.

\bibitem[Naidoo et~al.(2016)]{Naidoo2016}
K.~Naidoo, L.~Kazin, A.~Hamilton, et~al.
\newblock The Eridanus supervoid: confirmation, extent and photometric redshift
  constraints.
\newblock \emph{Monthly Notices of the Royal Astronomical Society},
  459(1):L71--L75, 2016.
\newblock \href{https://doi.org/10.1093/mnrasl/slw054}{doi:10.1093/mnrasl/slw054}.

\bibitem[Riess et~al.(2022)]{Riess2022}
A.~G.~Riess, W.~Yuan, L.~M.~Macri, et~al.
\newblock A comprehensive measurement of the local value of the Hubble constant
  with 1 km/s/Mpc uncertainty from the Hubble Space Telescope and the
  {SH0ES} team.
\newblock \emph{The Astrophysical Journal Letters}, 934(1):L7, 2022.
\newblock \href{https://doi.org/10.3847/2041-8213/ac5c5b}{doi:10.3847/2041-8213/ac5c5b}.

\bibitem[Kenworthy et~al.(2019)]{Kenworthy2019}
W.~D.~Kenworthy, D.~Scolnic, and A.~Riess.
\newblock The local perspective on the {Hubble} tension: local structure does
  not impact $H_0$.
\newblock \textit{The Astrophysical Journal}, 875(2):145, 2019.

\bibitem[Bose et~al.(2017)]{Bose2017}
S.~Bose, A.~Mazumdar, G.~W.~Morley, H.~Ulbricht, M.~Toro\v{s},
  M.~Paternostro, A.~A.~Geraci, P.~F.~Barker, M.~S.~Kim, and G.~Milburn.
\newblock Spin entanglement witness for quantum gravity.
\newblock \textit{Physical Review Letters}, 119:240401, 2017.

\bibitem[Marletto and Vedral(2017)]{Marletto2017}
C.~Marletto and V.~Vedral.
\newblock Gravitationally induced entanglement between two massive particles is
  sufficient evidence of quantum effects in gravity.
\newblock \textit{Physical Review Letters}, 119:240402, 2017.


\bibitem[Lovelock(1971)]{Lovelock1971}
D.~Lovelock.
\newblock The Einstein tensor and its generalizations.
\newblock \emph{Journal of Mathematical Physics}, 12(3):498--501, 1971.
\newblock \href{https://doi.org/10.1063/1.1665613}{doi:10.1063/1.1665613}.

\bibitem[Meiburg et~al.(2025)]{Meiburg2025LQI}
A.~Meiburg, L.~A.~Lessa, and R.~R.~Soldati.
\newblock A formalization of the generalized quantum {Stein's} lemma in {Lean}.
\newblock arXiv:2510.08672, 2025.
\newblock Lean-QuantumInfo:
  \href{https://github.com/Timeroot/Lean-QuantumInfo}{github.com/Timeroot/Lean-QuantumInfo}.

\bibitem[Sandeman(2026g)]{Sandeman2026RASM}
J.~F.~Sandeman.
\newblock Relational Actualism and the Standard Model ({RASM}).
\newblock Preprint, 2026.
\newblock \href{https://doi.org/10.5281/zenodo.19229335}{doi.org/10.5281/zenodo.19229335}.

\bibitem[Sandeman(2026RACL)]{Sandeman2026RACL}
J.~F.~Sandeman.
\newblock Relational Actualism: The Einstein Field Equations as the Unique
  Macroscopic Limit of the Actualization Graph ({RACL}).
\newblock \textit{Classical and Quantum Gravity} (submitted), 2026.
\newblock \href{https://doi.org/10.5281/zenodo.19270020}{doi.org/10.5281/zenodo.19270020}.

\bibitem[Sandeman(2026RATM)]{Sandeman2026RATM}
J.~F.~Sandeman.
\newblock Relational Actualism: The Topology of Matter ({RATM}).
\newblock \textit{Physical Review D} (submitted), 2026.
\newblock \href{https://doi.org/10.5281/zenodo.19265651}{doi.org/10.5281/zenodo.19265651}.

\bibitem[Sandeman(2026BDG)]{Sandeman2026BDG}
J.~F.~Sandeman.
\newblock Two Standard Model Coupling Constants from the Unique 4-Dimensional
  BDG Integers.
\newblock Preprint, 2026.
\newblock \href{https://doi.org/10.5281/zenodo.19324790}{doi.org/10.5281/zenodo.19324790}.

\bibitem[Sandeman(2026h)]{Sandeman2026Foundation}
J.~F.~Sandeman.
\newblock Relational Actualism: A Foundation Paper.
\newblock Preprint, 2026.
\newblock \href{https://doi.org/10.5281/zenodo.19229438}{doi.org/10.5281/zenodo.19229438}.

\bibitem[O'Hanlon(1972)]{OHanlon1972}
J.~O'Hanlon.
\newblock Intermediate-range gravity: a generally covariant model.
\newblock \textit{Physical Review Letters}, 29:137--138, 1972.

\bibitem[Damour and Esposito-Far\`{e}se(1992)]{Damour1992}
T.~Damour and G.~Esposito-Far\`{e}se.
\newblock Tensor-multi-scalar theories of gravitation.
\newblock \textit{Classical and Quantum Gravity}, 9:2093--2176, 1992.

\bibitem[Sotiriou and Faraoni(2010)]{Sotiriou2010}
T.~P.~Sotiriou and V.~Faraoni.
\newblock $f(R)$ theories of gravity.
\newblock \textit{Reviews of Modern Physics}, 82:451--497, 2010.

\end{thebibliography}

\end{document}
